% Options for packages loaded elsewhere
\PassOptionsToPackage{unicode}{hyperref}
\PassOptionsToPackage{hyphens}{url}
\PassOptionsToPackage{dvipsnames,svgnames,x11names}{xcolor}
%
\documentclass[
  10pt,
]{article}
\usepackage{amsmath,amssymb}
\usepackage{iftex}
\ifPDFTeX
  \usepackage[T1]{fontenc}
  \usepackage[utf8]{inputenc}
  \usepackage{textcomp} % provide euro and other symbols
\else % if luatex or xetex
  \usepackage{unicode-math} % this also loads fontspec
  \defaultfontfeatures{Scale=MatchLowercase}
  \defaultfontfeatures[\rmfamily]{Ligatures=TeX,Scale=1}
\fi
\usepackage{lmodern}
\ifPDFTeX\else
  % xetex/luatex font selection
  \setmainfont[]{DejaVu Serif}
  \setsansfont[]{DejaVu Sans}
  \setmonofont[]{DejaVu Sans Mono}
  \setmathfont[]{Latin Modern Math}
\fi
% Use upquote if available, for straight quotes in verbatim environments
\IfFileExists{upquote.sty}{\usepackage{upquote}}{}
\IfFileExists{microtype.sty}{% use microtype if available
  \usepackage[]{microtype}
  \UseMicrotypeSet[protrusion]{basicmath} % disable protrusion for tt fonts
}{}
\makeatletter
\@ifundefined{KOMAClassName}{% if non-KOMA class
  \IfFileExists{parskip.sty}{%
    \usepackage{parskip}
  }{% else
    \setlength{\parindent}{0pt}
    \setlength{\parskip}{6pt plus 2pt minus 1pt}}
}{% if KOMA class
  \KOMAoptions{parskip=half}}
\makeatother
\usepackage{xcolor}
\usepackage[margin=2.0cm]{geometry}
\usepackage{longtable,booktabs,array}
\usepackage{calc} % for calculating minipage widths
% Correct order of tables after \paragraph or \subparagraph
\usepackage{etoolbox}
\makeatletter
\patchcmd\longtable{\par}{\if@noskipsec\mbox{}\fi\par}{}{}
\makeatother
% Allow footnotes in longtable head/foot
\IfFileExists{footnotehyper.sty}{\usepackage{footnotehyper}}{\usepackage{footnote}}
\makesavenoteenv{longtable}
\setlength{\emergencystretch}{3em} % prevent overfull lines
\providecommand{\tightlist}{%
  \setlength{\itemsep}{0pt}\setlength{\parskip}{0pt}}
\setcounter{secnumdepth}{-\maxdimen} % remove section numbering
% doc_header.tex -- LaTeX preamble additions used by build_docs.sh for every
% Markdown -> LaTeX -> PDF conversion (the documents of the folder Revision_SM_Higgs).
\usepackage{etoolbox}
% longtable may not be loaded by pandoc when a document has no table
\usepackage{longtable}
% bold vectors: \bm -> \symbf under unicode-math (XeLaTeX), bm package otherwise
\makeatletter
\@ifpackageloaded{unicode-math}{\providecommand{\bm}[1]{\symbf{#1}}}{\usepackage{bm}}
\makeatother
% smaller type in long tables so that wide audit tables fit the text width;
% inside a table a line break is allowed after an underscore of plain text
% (pandoc writes it as \_), so that long identifiers wrap inside their column
% rows are set a little apart, so that symbols with indices do not touch
\AtBeginEnvironment{longtable}{\footnotesize\renewcommand{\arraystretch}{1.45}%
  \let\sptunderscore\_\def\_{\sptunderscore\allowbreak}}
\setlength{\LTcapwidth}{\textwidth}
% A short table is kept on one page.  The build (code/scripts/keep_tables.py) ends the rows of such a table with \\*,
% which forbids a page break after a row, and sets \sptkeptpenalty to 10000 in front of the table.  This penalty is
% put here in the place of the breakpoint (penalty 0) that longtable leaves after every table, so that the paragraph
% that follows a kept table, its caption, begins on the same page; it is set back to 0 after the table.
\newcount\sptkeptpenalty
\sptkeptpenalty=0
\makeatletter
\patchcmd{\endlongtable}{\endgraf\penalty\z@\addvspace\LTpost}%
  {\endgraf\penalty\sptkeptpenalty\addvspace\LTpost\global\sptkeptpenalty=\z@}%
  {}{\PackageError{doc_header}{the end of longtable could not be patched}{The tables can still be typeset, but a short table may be separated from its caption.}}
\makeatother
% allow line breaks after slashes and hyphens inside table cells
\usepackage{xurl}
% code blocks: break over-long lines instead of running into the margin.
% Plain ``` blocks are typeset by pandoc as {verbatim}, highlighted ones
% (```bash) as {Highlighting}; both are re-declared as breakable fancyvrb
% environments.  The continuation line is marked with a small hooked arrow.
\usepackage{fvextra}
\RecustomVerbatimEnvironment{verbatim}{Verbatim}{breaklines,breakanywhere,fontsize=\small}
\makeatletter
\@ifundefined{Highlighting}{}{%
  \RecustomVerbatimEnvironment{Highlighting}{Verbatim}{breaklines,breakanywhere,fontsize=\small,commandchars=\\\{\}}}
\makeatother
% inline code (pandoc writes it as \texttt{...}): never hyphenate -- an inserted
% hyphen could be read as part of a file name -- but allow a line break after
% an underscore, slash, comma, colon or equals sign, so that long paths and
% column lists wrap inside the text width and inside table columns.
% Inside a formula (e.g. the name of a notebook symbol in an equation) the same
% command must first leave math mode, which \text of amsmath does.
\ExplSyntaxOn
\tl_new:N \l__spt_code_tl
\cs_new_protected:Npn \__spt_texttt:n #1
  {
    \tl_set:Nn \l__spt_code_tl {#1}
    \regex_replace_all:nnN { ([\/\,\:\=]) } { \1 \c{allowbreak} } \l__spt_code_tl
    \regex_replace_all:nnN { \c{_} } { \c{_} \c{allowbreak} } \l__spt_code_tl
    \group_begin: \ttfamily \hyphenchar \font = -1 \scan_stop: \l__spt_code_tl \group_end:
  }
\RenewDocumentCommand \texttt { m }
  {
    \mode_if_math:TF { \text { \__spt_texttt:n {#1} } } { \__spt_texttt:n {#1} }
  }
\ExplSyntaxOff
% A figure is written in the Markdown source as an image paragraph followed by
% a caption paragraph ("**Figure S1.** ..."), so that the caption is also
% visible on GitHub.  An image that starts a paragraph is centered.
\usepackage{graphicx}% pandoc loads it only for documents that contain an image
\NewCommandCopy\sptincludegraphics\includegraphics
\RenewDocumentCommand\includegraphics{O{} m}{%
  \ifvmode\begin{center}\sptincludegraphics[#1]{#2}\end{center}%
  \else\sptincludegraphics[#1]{#2}\fi}
% keep figures where they are referenced
\usepackage{float}
\floatplacement{figure}{H}
% running header with the document title
\usepackage{fancyhdr}
\pagestyle{fancy}
\fancyhf{}
\fancyhead[L]{\small PROJECT spacetimes, Revision\_SM\_Higgs}
\fancyhead[R]{\small\thepage}
\renewcommand{\headrulewidth}{0.4pt}
% 16x16 matrices: tighter columns so that a full matrix fits the text width
\setlength{\arraycolsep}{3.2pt}
% The documents of this project contain long tagged symbols, file paths and
% identifiers that cannot be hyphenated.  A paragraph is therefore allowed to
% have looser lines before TeX gives up and lets a line stick out of the margin.
\AtBeginDocument{\tolerance=2000 \emergencystretch=5em \hbadness=2000}
\ifLuaTeX
  \usepackage{selnolig}  % disable illegal ligatures
\fi
\IfFileExists{bookmark.sty}{\usepackage{bookmark}}{\usepackage{hyperref}}
\IfFileExists{xurl.sty}{\usepackage{xurl}}{} % add URL line breaks if available
\urlstyle{same}
\hypersetup{
  pdftitle={The provenance of the Mathematica notebook superfluid.nb},
  pdfauthor={Prepared for Patrick Lee Nash, Ph.D.},
  colorlinks=true,
  linkcolor={blue},
  filecolor={Maroon},
  citecolor={blue},
  urlcolor={blue},
  pdfcreator={LaTeX via pandoc}}

\title{The provenance of the Mathematica notebook superfluid.nb}
\usepackage{etoolbox}
\makeatletter
\providecommand{\subtitle}[1]{% add subtitle to \maketitle
  \apptocmd{\@title}{\par {\large #1 \par}}{}{}
}
\makeatother
\subtitle{Provenance document \texttt{P3} of PROJECT spacetimes, folder
Revision\_SM\_Higgs}
\author{Prepared for Patrick Lee Nash, Ph.D.}
\date{November 21, 2026}

\begin{document}
\maketitle

\textbf{The state of this document under the revised hypotheses (21 and
22 November 2026).} Sections 1 to 9 were written for the notebook of the
earlier hypotheses, under which the spinor field was the DiracLike
spinor with the electric charge 0 and
\(SU(2)_{\rm v}\times U(1)^{\rm F}_{\rm LL}\) was a new gauge group;
Parts I to III of the notebook state those hypotheses. Under the revised
hypotheses of this folder, in their second version, \(SU(2)_{\rm v}\) is
\(SU(2)_L\), the isodoublet \(\Psi_{\rm LL}\) is joined by the singlets
\(\chi_{\rm SH}\) (hypercharge \(\tfrac14\)) and \(\chi'_{\rm SH}\)
(hypercharge \(-\tfrac34\)), and the superfluid is the condensate of the
pairs of a particle of \(\Psi_{\rm LL}\) with a particle of a singlet,
whose two pair fields have the quantum numbers of the Higgs doublet.
Part R of the notebook, added on 22 November 2026, states the revised
hypotheses and repeats the checks of the four stages of the revision;
Section 10 records it, and Sections 2 to 6 and 9 name it where they
describe the whole notebook.

\hypertarget{what-this-document-is}{%
\section{1. What this document is}\label{what-this-document-is}}

This document records how the Mathematica notebook
\texttt{superfluid.nb} of the folder \texttt{Revision\_SM\_Higgs} was
written, built and evaluated, what it computes, which checks it makes
and how each check is repeated by a program in Rust and a program in
Python, and which errors were found and corrected while it was written.
Every statement about the results is taken from the log of the
evaluation of the notebook in the front end of Wolfram (the file
\texttt{data/hck\_notebook\_log.txt}) and from the reports of the Rust
and the Python programs, by the program
\texttt{code/scripts/write\_notebook\_fragments.py}; no number of the
tables below is typed by hand.

\hypertarget{the-task-and-the-state}{%
\section{2. The task and the state}\label{the-task-and-the-state}}

\textbf{The task} (item 5 of the main task, in my words). A new,
stand-alone Mathematica notebook \texttt{superfluid.nb}, written in the
Wolfram language only, that evaluates without errors or warnings and is
an encyclopedia of the whole solution: every symbol, documented by
comments; every equation that can be executed; every solution; the
explanations, the figures, the graphs and the videos, each documented.
For every executable part of the notebook, executable pure Rust code and
Python code are part of the solution, with complete instructions, so
that users without a Wolfram license can follow; every time integration
in the notebook uses the method IDA inside the function \texttt{NDSolve}
(never a simple fixed-step integrator such as the fourth-order
Runge-Kutta method).

\textbf{The state of this version.} The notebook is written chapter by
chapter, and each chapter is committed as soon as it is complete. This
version contains the introduction, the register of the checks (chapter
0), Part I (items 1 to 5 of the problem: the conventions, the manifold
M31 and the Lorentz super-loop, the spinor field and theorem D1, the
gauge group GEE and its covariant derivatives, the Lagrangian and the
field equations, and the four analytic solutions), Part II (item 6 of
the problem: the pair channels and the gauge bosons that mediate the
pairing, the three methods and their numerical engine, the recursion of
the potential to \(V^{(2)}_{\rm SF}\), the ground states and their gaps,
the Goldstone and the collective modes, the susceptibilities and the
masses of the gauge bosons, the exchange of the gauge bosons), the first
section of Part III (the Landau form of the effective potentials: its
definition with examples, its accuracy, the cause of its error and why
the error matters; the answer to the author's questions of 21 November
2026, which he asked to be kept in the notebook), and the second section
of Part III (Section III.T: the free energy of the five methods at
finite temperature, the critical temperatures, the gap equation, the
size of the order parameter against the radius of convergence, and the
remark that the free energies of the methods are mean-field free
energies, with its importance; the addendum asked for by the author on
21 November 2026), and the third section of Part III (Section III.G: the
gradient coefficient of the field at the critical temperature from the
static pair response of a moving pair, the two regularizations, the
slope and the quartic coefficient, the Ginzburg number and the width of
the critical region; the superfluid as the Higgs sector of its own gauge
group, and why it cannot be the Higgs field of the Standard Model (the
DiracLike spinor has electric charge 0 and no electroweak quantum
numbers), the couplings of the amplitude field by exact differentiation,
and whether the Higgs is a composite particle; the second addendum asked
for by the author on 21 November 2026, with its three remarks; and, at
the end of Part II, the remark on the zeugon, the massless gauge boson
of the superfluid), with 181 checks. The rest of Part III and Part IV
(the effective Lagrangians and the Mexican hat, the masses of the gauge
bosons in every vacuum, the surviving symmetry and the dynamics with
IDA) follow in the next versions. The first section of Part III takes
the five computed polar curves of the effective potential from the Rust
program \texttt{sf\_part3} (a closed cell of data with the SHA-256
checksums of its files) and recomputes every number from them; Section
III.T computes the thermal states, the critical temperatures and the
Taylor coefficients itself with the engine of Part II, and takes the
free-energy curves of the Rust program from a third closed cell; Section
III.G computes the static pair response, the coefficients and the
couplings itself with the engines of Part II and of Section III.T and
takes the results of the modules \texttt{ginzburg} and \texttt{higgs} of
the Rust program from a fourth closed cell; Part II computes all its
states itself and compares them with the values of the Rust program
\texttt{sf\_part2} (another closed cell of data). No part so far has a
time integration.

\textbf{Part R (22 November 2026).} The notebook now ends with Part R,
the revised hypotheses (Section 10): the quantum numbers of the three
fields and of their pairs (stage 2 of the revision), the ground state of
the superfluid (stage 3), the masses of the gauge bosons and the modes
(stage 4) and the dynamics in time (stage 6), with the remarks R8 to
R12, in 28 checks; with them the notebook has 181 checks. Part R takes
the tables of the Rust programs of the revision from a fifth closed cell
of data and integrates the reduced model of stage 6 with the method IDA
inside \texttt{NDSolve}.

\hypertarget{the-files}{%
\section{3. The files}\label{the-files}}

The files of the notebook, each with its SHA-256 checksum:

\begin{itemize}
\tightlist
\item
  \texttt{superfluid.nb}: the notebook, evaluated in the front end of
  Wolfram (inputs, printed checks, the table of all checks);\\
  \texttt{133a3b14343b90f2adbe17d31d038639f2bb53feab1b2f0a5ffe9afabe570b99}
\item
  \texttt{superfluid\_notebook.pdf}: the evaluated notebook printed by
  the front end, for readers without Wolfram;\\
  \texttt{13b2f9bfae8ede21eb4afb8e055ab6e53d3c2586268f42f7b846d9cbf76ec2ee}
\item
  \texttt{wolfram/source/superfluid.source.wl}: the source of the
  notebook: the title, the introduction, chapter 0 (the register of the
  checks) and the final list of all checks; it includes the sources of
  the parts;\\
  \texttt{182d76d75e5416c249b6979690253b441c7078ef329d6f88e3018bc15b789e5c}
\item
  \texttt{wolfram/source/part1.source.wl}: the source of Part I;\\
  \texttt{82e6e590de5995275a920f1a96ea97c43d49ab364d02a7a01988b81110eb0879}
\item
  \texttt{wolfram/source/part2.source.wl}: the source of Part II;\\
  \texttt{c27a6c74bbce09f66dace6df4c7732bbfde41ab1a1eb0301b231dcca3f0e01d9}
\item
  \texttt{wolfram/source/data\_part2.wl}: the values of Part II of the
  Rust program and its table of pair-channel energies, included in a
  closed cell of the notebook (written by
  \texttt{wolfram/tools/make\_data\_cells.py}; its first lines name the
  data files and their SHA-256 checksums);\\
  \texttt{31db4e03444eeb963786e49252e7078316e8fe18914074db910696e938e49c30}
\item
  \texttt{wolfram/source/part3\_landau.source.wl}: the source of the
  first section of Part III (the Landau form);\\
  \texttt{9cc6176d7163c81c26b7d67057df2e055ad6448aa31781174610e8ddd63d2c9d}
\item
  \texttt{wolfram/source/data\_part3.wl}: the values of Part III of the
  Rust program and its five polar curves, included in a closed cell of
  the notebook (written by \texttt{wolfram/tools/make\_data\_cells.py};
  its first lines name the data files and their SHA-256 checksums);\\
  \texttt{8a124c8a2cf1382aa9189b5c2a8d210735cc3aeb19d22afb0bece0d3e184116a}
\item
  \texttt{wolfram/source/part3\_thermal.source.wl}: the source of the
  second section of Part III (Section III.T, the free energy at finite
  temperature);\\
  \texttt{673d6b22ba4177fe636a717d664bcef3d7c0f6b612cbe21a868ad726c4835caf}
\item
  \texttt{wolfram/source/data\_part3\_thermal.wl}: the thermal results
  of the Rust program (the table of the order parameter, the residuals
  of the Landau forms and the free-energy curves of the five methods),
  included in a closed cell of the notebook (written by
  \texttt{wolfram/tools/make\_data\_cells.py\ 3T}; its first lines name
  the data files and their SHA-256 checksums);\\
  \texttt{5f3fe6ec541cd7e8a13c08517a63ba8b10deff9a5b37eded75aec6bbef2ebfac}
\item
  \texttt{wolfram/source/part3\_ginzburg.source.wl}: the source of the
  third section of Part III (Section III.G, the width of the critical
  region and the superfluid as a Higgs sector);\\
  \texttt{ce3d808ecca2c26e66901435fd14df314dd1c6c308288810e925b627e34b0ad8}
\item
  \texttt{wolfram/source/data\_part3\_ginzburg.wl}: the results of the
  modules \texttt{ginzburg} and \texttt{higgs} of the Rust program (the
  coefficients of the five methods, the two regularizations and the
  numbers of the Higgs sector, one row per method), included in a closed
  cell of the notebook (written by
  \texttt{wolfram/tools/make\_data\_cells.py\ 3G}; its first lines name
  the data files and their SHA-256 checksums);\\
  \texttt{8aabc404d7c5ef67f5670eb73da39446c8289e240d299706d9b409fc6344089c}
\item
  \texttt{wolfram/source/part\_revision.source.wl}: the source of Part R
  (the revised hypotheses; Section 10);\\
  \texttt{16b1307c9db03357369c791a8d2a92b401e96ca9dc5c4af0a2d7d7cb02c9f8bf}
\item
  \texttt{wolfram/source/data\_partR.wl}: the tables of the Rust
  programs of the revision (the channels of stage 2, the normal vacuum
  and the ground states of stage 3, the minimum, the masses and the
  values of stage 4, the runs and, of the signals of stage 6, the rows
  that the checks use), every field as a string, included in a closed
  cell of the notebook (written by
  \texttt{wolfram/tools/make\_data\_cells.py\ R}; its first lines name
  the data files and their SHA-256 checksums);\\
  \texttt{0df097f31f5630e03838f8054ff6c696e278b693383ffb571dfa08a5ad4a50c2}
\item
  \texttt{wolfram/tools/make\_data\_cells.py}: writes the data cells of
  the notebook from the data files of the Rust programs;\\
  \texttt{89584b3fb5c1bd6d576f076695be37ec716b88aa0157b6d5acf497354738fcbb}
\item
  \texttt{wolfram/tools/build\_notebook.wls}: the builder: turns the
  source into the notebook (text cells with formulas, input cells parsed
  by the front end);\\
  \texttt{ff661ed5051e22e8d14487cb80f91fb9cbedb2243bd550d6e6c688afb83fd94a}
\item
  \texttt{wolfram/tools/sf\_formula.wl}: the reader of formulas: turns
  the TeX of the formulas of the texts into two-dimensional boxes;\\
  \texttt{686f34fb8f6e985cb7ea376cc41050b6d53ad2035a8acca00f39c30490575466}
\item
  \texttt{wolfram/tools/guirun.wl}: the evaluation of the notebook in
  the front end, in a fresh kernel, with the log;\\
  \texttt{9492bd3e8317cccc365e8b3f04168cfc7bb0cb1db8f5d94ee1d0d18fcc219173}
\item
  \texttt{wolfram/tools/gui\_start.wl}: one of the three calls of the
  evaluation (Section 4);\\
  \texttt{8a5a9b4be3ca32a81bd62461e8f353c7b8ee22e1a9b9a96770380060ede83e5f}
\item
  \texttt{wolfram/tools/gui\_go.wl}: one of the three calls of the
  evaluation (Section 4);\\
  \texttt{1ad800de5df7a97cc38ed4755ce8f3769bce4c59adf2ed62bdb6c3b09a4e5a77}
\item
  \texttt{wolfram/tools/gui\_status.wl}: one of the three calls of the
  evaluation (Section 4);\\
  \texttt{0bc716d56ccc71443630e2d77ba6bfb00430660c80db499503a29bbde046967c}
\item
  \texttt{wolfram/tools/make\_dryrun.py}: writes the code of all input
  cells into one file for a test run by \texttt{wolframscript} without
  the front end;\\
  \texttt{e08e498dbae36ade3664a5a4a4ffdac4da5ee9ecc45a1189f25e03d18eca2fa1}
\item
  \texttt{data/hck\_notebook\_log.txt}: the log of the evaluation: the
  version of Wolfram, the counts of the cells, every message and printed
  line, one line per check;\\
  \texttt{8599108efc70155d07a6c7d93cdbdc9729b2803640afdc558247dc60e1f9b29e}
\item
  \texttt{data/hck\_notebook\_results.wl}: the register of the checks,
  written by the kernel at the end of the evaluation;\\
  \texttt{3a785c6515d231be2f49b909ec0c2d505558b53eca87df8a1dd209b812ce7a93}
\item
  \texttt{data/hck\_notebook\_build.txt}: the report of the builder when
  it built the notebook that was evaluated (the counts of the cells, of
  the closed cells of data, of the formulas in the texts and of the
  warnings);\\
  \texttt{c2a48bd1d3cb7f96d0122e0edff448bab2f5aa1f3e6e32fc3a44ab685b8a9610}
\item
  \texttt{code/scripts/write\_notebook\_fragments.py}: writes the tables
  and numbers of this document from the log and the reports;\\
  \texttt{732101dc280ca1a37ccf6187a2f544143ca60a2260b26bb7f3093007de028ecd}
\end{itemize}

The builder, the reader of formulas and the evaluation script were first
written for the notebook \texttt{3+1super-loop.nb} of the project; they
are copied into this folder with their contexts renamed (from
\texttt{SL31} to \texttt{SFNB}), so that this folder holds every program
that makes its notebook. The reader of formulas there took its table of
the characters of TeX commands from a second package (the reader of
symbols of that notebook), which this notebook does not use; the table
is now part of \texttt{sf\_formula.wl} itself (error E2 of Section 9).

\hypertarget{how-the-notebook-is-built-and-evaluated}{%
\section{4. How the notebook is built and
evaluated}\label{how-the-notebook-is-built-and-evaluated}}

\textbf{For a reader with Wolfram.} Open \texttt{superfluid.nb} in
Wolfram (version 15 or later; this version was evaluated with 15.0.1 for
Mac OS X ARM (64-bit) (July 2, 2026)) and choose Evaluation
\textgreater{} Evaluate Notebook. Nothing outside the notebook is read.

\textbf{How the notebook is made from its source.} The source is a text
in the package format of the Wolfram language, in which every cell
begins with a marker line such as \texttt{(*\ ::Text::\ *)} or
\texttt{(*\ ::Input::\ *)}. On the author's computer:

\begin{verbatim}
wolframscript -file wolfram/tools/build_notebook.wls wolfram/source/superfluid.source.wl in/superfluid.nb \
    wolfram/tools/sf_formula.wl
\end{verbatim}

The builder replaces the lines
\texttt{(*\ @@INCLUDECELLS\ part1.source.wl@@\ *)},
\texttt{(*\ @@INCLUDECELLS\ part2.source.wl@@\ *)},
\texttt{(*\ @@INCLUDECELLS\ part3\_landau.source.wl@@\ *)},
\texttt{(*\ @@INCLUDECELLS\ part3\_thermal.source.wl@@\ *)},
\texttt{(*\ @@INCLUDECELLS\ part3\_ginzburg.source.wl@@\ *)} and
\texttt{(*\ @@INCLUDECELLS\ part\_revision.source.wl@@\ *)} of the main
source by the cells of these parts, a line such as
\texttt{(*\ @@INCLUDE\ data\_part3.wl@@\ *)} inside an input cell by
that file, turns every formula of a text (TeX between dollar signs) into
boxes with the reader of formulas, lets the front end parse every input
cell (so that comments, line breaks and indentation stay as written),
and reports every formula that the reader cannot read as a warning. This
version: 403 cells (5 of them closed cells of data), 1574 formulas in
the texts, 0 warnings (the builder's own count, printed when it ran and
kept in \texttt{data/hck\_notebook\_build.txt}).

\textbf{How it is evaluated.} The notebook is evaluated in the front end
of the running Wolfram application, exactly as by the menu command
Evaluate Notebook, but in a fresh kernel named \texttt{SFNBRun}, in
three calls (each returns at once):

\begin{enumerate}
\def\labelenumi{\arabic{enumi}.}
\tightlist
\item
  \texttt{gui\_start.wl}: opens the built notebook, sets its kernel, and
  starts the kernel with a hidden cell that makes the kernel append the
  number of every input line and the time to a progress file;
\item
  \texttt{gui\_go.wl}: evaluates all cells, and after them a hidden cell
  that writes the register of the checks
  (\texttt{SFNB}Test\texttt{\$Results}) to the results file and quits
  the kernel;
\item
  \texttt{gui\_status.wl}: saves the evaluated notebook, prints it to
  PDF, and writes the log: the version, the counts of the cells, every
  message and printed line, and one line per check.
\end{enumerate}

On the author's computer these calls are made by a small program that
runs inside the Wolfram application and executes requests left as files
in a work folder (\texttt{.git-transfer/mma/} of the author's copy of
the repository, which is not part of the repository); its handler
accepts a request \texttt{wls\ \textless{}script\textgreater{}} (run by
\texttt{wolframscript} in a fresh kernel) or
\texttt{gui\ \textless{}file\textgreater{}} (evaluated in the kernel of
the application). The first evaluation of this version made the three
calls with the same two Wolfram lines that the files
\texttt{gui\_start.wl}, \texttt{gui\_go.wl} and \texttt{gui\_status.wl}
now hold with comments.

\textbf{The dry run.} Before every build,
\texttt{python3\ wolfram/tools/make\_dryrun.py\ OUT.wls\ wolfram/source/superfluid.source.wl}
writes the code of all input cells, in order, into one file, and
\texttt{wolframscript\ -file\ OUT.wls} evaluates it in a plain kernel
(243 seconds in this version, of which Part II takes about two minutes,
Section III.T about half a minute, Section III.G about 80 seconds and
Part R about 3 seconds, by the times that the dry run prints after its
checks); the only difference from the notebook is that it prints the
computing time after every check. The dry run finds errors of the code
in seconds; the evaluation in the front end stays the only evidence for
the results reported here.

\hypertarget{the-evaluation-of-this-version}{%
\section{5. The evaluation of this
version}\label{the-evaluation-of-this-version}}

15.0.1 for Mac OS X ARM (64-bit) (July 2, 2026); 305 seconds; 93 input
cells, 93 of them evaluated; 16 output cells; 211 printed lines;
\textbf{0 messages} (no error and no warning of the Wolfram language);
verdict: \textbf{ALL PASS (181 of 181), 0 messages}.

\hypertarget{the-checks-and-their-counterparts}{%
\section{6. The checks and their
counterparts}\label{the-checks-and-their-counterparts}}

Every check of the notebook carries the identifier of the same check in
the Rust program with the prefix W, and names its counterparts in the
Rust program (\texttt{rust/sfpart1}, program \texttt{sf\_part1}, for
Part I; \texttt{rust/sfpart2}, program \texttt{sf\_part2}, for Part II;
\texttt{rust/sfpart3}, program \texttt{sf\_part3}, for Part III, whose
module \texttt{thermal} makes the checks P3-T1 to P3-T9 of Section III.T
and whose modules \texttt{ginzburg} and \texttt{higgs} make the checks
P3-G1 to P3-G12 and P3-H1 to P3-H3 of Section III.G) and in the Python
program (\texttt{python/sf\_part1.py}, \texttt{python/sf\_part2.py},
\texttt{python/sf\_part3.py}). Where the Rust program checks a statement
in two steps, the notebook names both checks (``P2-6b, P2-6c''); each
must have passed. The checks of Part R carry the identifiers of the
Wolfram scripts of the revision instead: W, the letters of the stage
(QN, SF, H, D) and a number (decision D34 of the task record). They name
as counterparts the checks of the Rust programs \texttt{sh\_quantum},
\texttt{sh\_superfluid}, \texttt{sh\_higgs} and \texttt{sh\_dynamics}
and of the Python programs \texttt{python/sh\_quantum.py},
\texttt{sh\_superfluid.py}, \texttt{sh\_higgs.py} and
\texttt{sh\_dynamics.py} that test the same statements, as ranges
(``QN-1a to QN-1e'', each of which must have passed) or ``none'' where
no program of the other language tests the statement. The table gives,
for each check, the deviation measured by the notebook and by each
counterpart (as their reports state it; ``--'' for a check that is
decided by facts, such as a count or an exact identity). The program
that writes the table stops with an error if a counterpart is missing
from its report or did not pass.

\begin{longtable}[]{@{}
  >{\raggedright\arraybackslash}p{(\columnwidth - 12\tabcolsep) * \real{0.0825}}
  >{\raggedright\arraybackslash}p{(\columnwidth - 12\tabcolsep) * \real{0.0619}}
  >{\raggedright\arraybackslash}p{(\columnwidth - 12\tabcolsep) * \real{0.1134}}
  >{\raggedright\arraybackslash}p{(\columnwidth - 12\tabcolsep) * \real{0.1031}}
  >{\raggedright\arraybackslash}p{(\columnwidth - 12\tabcolsep) * \real{0.1340}}
  >{\raggedright\arraybackslash}p{(\columnwidth - 12\tabcolsep) * \real{0.1340}}
  >{\raggedright\arraybackslash}p{(\columnwidth - 12\tabcolsep) * \real{0.3711}}@{}}
\toprule\noalign{}
\begin{minipage}[b]{\linewidth}\raggedright
check
\end{minipage} & \begin{minipage}[b]{\linewidth}\raggedright
result
\end{minipage} & \begin{minipage}[b]{\linewidth}\raggedright
deviation in the notebook
\end{minipage} & \begin{minipage}[b]{\linewidth}\raggedright
tolerance
\end{minipage} & \begin{minipage}[b]{\linewidth}\raggedright
Rust: check, deviation
\end{minipage} & \begin{minipage}[b]{\linewidth}\raggedright
Python: check, deviation
\end{minipage} & \begin{minipage}[b]{\linewidth}\raggedright
statement
\end{minipage} \\
\midrule\noalign{}
\endhead
\bottomrule\noalign{}
\endlastfoot
WSM-1 & PASS & -- & -- & SM-1 (PASS), \(4.163\times10^{-16}\) & SM-1
(PASS), \(3.863\times10^{-16}\) & Standard Model with +i g: the
eigenvalues of the mass matrix are 0, M\_W\^{}2 (twice), M\_Z\^{}2, and
the photon sin(thetaW) W\^{}3 + cos(thetaW) B is massless {[}exact, with
symbolic g, g', v{]} \\
WA-2 & PASS & \(0\) & \(0\) & A-2 (PASS), \(0\) & A-2 (PASS), \(0\) &
bi-octonion units: e\_p\^{}2 = -1, e\_p e\_q = -e\_q e\_p, e\_j e\_4 =
e\_(4+j) (exact) \\
WA-3 & PASS & \(0\) & \(0\) & A-3 (PASS), \(0\) & A-3 (PASS), \(0\) &
N(xy) = N(x) N(y), 200 random Gaussian-integer bi-octonions (exact) \\
WA-4 & PASS & \(0\) & \(0\) & A-4 (PASS), \(0\) & A-4 (PASS), \(0\) & x
xbar = N(x) e\_0 (exact) \\
WA-5 & PASS & \(0\) & \(0\) & A-5 (PASS), \(0\) & A-5 (PASS), \(0\) &
alternativity (xx)y = x(xy), (yx)x = y(xx) (exact) \\
WA-6 & PASS & \(0\) & \(0\) & A-6 (PASS), \(0\) & A-6 (PASS), \(0\) &
the three Moufang identities (exact) \\
WA-7 & PASS & -- & -- & A-7 (PASS), -- & A-7 (PASS), -- & 168 of the 343
triples of imaginary units have a nonzero associator; (e1 e2) e4 - e1
(e2 e4) = 2 e7 {[}168 nonzero{]} \\
WA-8 & PASS & \(0\) & \(0\) & A-8 (PASS), \(1.776\times10^{-15}\) & A-8
(PASS), \(7.589\times10^{-15}\) & M(bb') = M(b)M(b'), det M = N,
M(b\^{}dag) = M(b)\^{}dag, inverse map (exact) \\
WA-10b & PASS & \(4.885\times10^{-15}\) & \(1.000\times10^{-10}\) &
A-10b (PASS), \(5.551\times10^{-15}\) & A-10b (PASS),
\(5.329\times10^{-15}\) & L\^{}T eta L = eta and det L = 1 (200 random
unit g) \\
WA-10c & PASS & -- & -- & A-10c (PASS), -- & A-10c (PASS), -- & L\_00
\textgreater= 1: orthochronous {[}200 random unit g{]} \\
WA-10d & PASS & \(5.329\times10^{-15}\) & \(1.000\times10^{-10}\) &
A-10d (PASS), \(4.441\times10^{-15}\) & A-10d (PASS),
\(7.105\times10^{-15}\) & L(gh) = L(g) L(h) and L(-g) = L(g) (200 random
pairs) \\
WA-12 & PASS & \(0\) & \(0\) & A-12 (PASS), \(0\) & A-12 (PASS), \(0\) &
spinor charges under the super-loop multiplications (exact) \\
WF-2 & PASS & \(0\) & \(0\) & F-2 (PASS), \(0\) & F-2 (PASS), \(0\) &
Clifford algebra, chirality, charge conjugation (exact) \\
WF-3a & PASS & \(1.777\times10^{-14}\) & \(1.000\times10^{-10}\) & F-3a
(PASS), \(1.599\times10^{-14}\) & F-3a (PASS), \(1.604\times10^{-14}\) &
S\^{}-1 gamma\^{}p S = sum\_q L\_pq gamma\^{}q (200 random g) \\
WF-3b & PASS & \(8.951\times10^{-16}\) & \(1.000\times10^{-10}\) & F-3b
(PASS), \(6.280\times10^{-16}\) & F-3b (PASS), \(7.448\times10^{-16}\) &
S\^{}dag gamma\^{}0 S = gamma\^{}0 (200 random g) \\
WF-3d & PASS & \(3.762\times10^{-15}\) & \(1.000\times10^{-10}\) & F-3d
(PASS), \(3.625\times10^{-15}\) & F-3d (PASS), \(3.827\times10^{-15}\) &
the Dirac matrices of the transformed frame are S gamma\^{}p S\^{}-1
(200 random g) \\
WF-5a & PASS & \(0\) & \(0\) & F-5a (PASS), \(0\) & F-5a (PASS), \(0\) &
D1: the Pauli-Guersey generators commute with the Lorentz generators
(exact) \\
WF-5c & PASS & -- & -- & F-5c (PASS), -- & F-5c (PASS), -- & D1: the
Lorentz-invariant bilinears of one point form a space of dimension 6;
those invariant also under the Pauli-Guersey SU(2), dimension 0
{[}dimensions 6, 0 (exact ranks){]} \\
WC-1 & PASS & \(5.551\times10^{-16}\) & \(1.000\times10^{-9}\) & C-1
(PASS), \(5.551\times10^{-16}\) & C-1 (PASS), \(3.736\times10^{-14}\) &
SU(2)\_v x U(1)\^{}F covariance of nabla on Psi (symbolic derivatives, 4
points) \\
WC-3 & PASS & \(1.178\times10^{-16}\) & \(1.000\times10^{-8}\) & C-3
(PASS), \(2.238\times10^{-16}\) & C-3 (PASS), \(2.755\times10^{-13}\) &
field strength W'\_pq = U W\_pq U\^{}dag (symbolic second derivatives, 2
points) \\
WC-4 & PASS & \(5.881\times10^{-16}\) & \(1.000\times10^{-9}\) & C-4
(PASS), \(9.037\times10^{-16}\) & C-4 (PASS), \(6.969\times10^{-14}\) &
isotriplet: nabla Phi = dPhi - g W x Phi + 2 i g\_F B Phi is covariant
(symbolic derivatives) \\
WC-8 & PASS & \(1.256\times10^{-15}\) & \(1.000\times10^{-9}\) & C-8
(PASS), \(2.809\times10^{-15}\) & C-8 (PASS), \(5.043\times10^{-14}\) &
local SL(2,C): gamma\^{}p{[}f'{]} (d\_p + Gamma'\_p)(S psi) = S
gamma\^{}p d\_p psi (symbolic derivatives) \\
WE-2 & PASS & \(0\) & \(0\) & E-2 (PASS), \(0\) & E-2 (PASS), \(0\) &
scalar contact term: {[}Psi\_e, H\_S{]} = G\^{}S (Psi-bar
Psi)(gamma\^{}0 Psi)\_e (exact, 256 states) \\
WE-3 & PASS & \(0\) & \(0\) & E-3 (PASS), \(3.553\times10^{-15}\) & E-3
(PASS), \(3.553\times10^{-15}\) & vector contact term (exact, 256
states) \\
WE-4 & PASS & \(0\) & \(0\) & E-4 (PASS), \(0\) & E-4 (PASS), \(0\) &
pairing term: {[}Psi\_e, H\_P{]} = -(G/2) sum (K\^{}dag Psi\^{}*)\_e Pi
(exact, 256 states) \\
WE-14 & PASS & -- & -- & E-14 (PASS), \(2.331\times10^{-15}\) & E-14
(PASS), -- & symbolic Euler-Lagrange equations of the Yang-Mills
Lagrangian with the source: d\_p W\^{}(a,p p') - g eps\_abc W\^{}b\_p
W\^{}(c,p p') = g j\^{}(a,p') (12 components) {[}12 of 12 equations
identical (EulerEquations; signs \{1, 1, 1, 1, 1, 1, 1, 1, 1, 1, 1,
1\}){]} \\
WAS1-3 & PASS & \(1.776\times10^{-15}\) & \(1.000\times10^{-12}\) &
AS1-3 (PASS), \(2.887\times10^{-15}\) & AS1-3 (PASS),
\(1.554\times10^{-15}\) & spin-isospin locked field: four analytic
branches (61 momenta) \\
WAS2-2 & PASS & \(4.828\times10^{-3}\) & \(1.000\times10^{-2}\) & AS2-2
(PASS), \(4.828\times10^{-3}\) & AS2-2 (PASS), \(4.828\times10^{-3}\) &
Landau levels, Wilson lattice spacing 0.025 (largest relative deviation
of the seven levels) \\
WAS2-4 & PASS & \(6.450\times10^{-5}\) & \(1.000\times10^{-3}\) & AS2-4
(PASS), \(6.450\times10^{-5}\) & AS2-4 (PASS), \(6.450\times10^{-5}\) &
Landau levels, Richardson extrapolation 2 E(0.025) - E(0.05) \\
WAS3-1 & PASS & \(2.665\times10^{-15}\) & \(1.000\times10^{-12}\) &
AS3-1 (PASS), \(3.775\times10^{-15}\) & AS3-1 (PASS),
\(2.665\times10^{-15}\) & BdG spectrum of unitary pair fields (50
random) \\
WAS3-4 & PASS & \(3.005\times10^{-11}\) & \(1.000\times10^{-9}\) & AS3-4
(PASS), -- & AS3-4 (PASS), \(3.005\times10^{-11}\) & the gap at G = 22.5
agrees with the Rust value 0.2094933033 \\
WAS4-5a & PASS & -- & -- & AS4-5a (PASS), -- & AS4-5a (PASS), -- & three
bound states below the continuum for each spacing {[}counts \{3, 3,
3\}{]} \\
WAS4-5b & PASS & \(1.116\times10^{-5}\) & \(1.000\times10^{-3}\) &
AS4-5b (PASS), \(1.116\times10^{-5}\) & AS4-5b (PASS),
\(1.116\times10^{-5}\) & kink bound states, Richardson extrapolation \\
WAS4-3 & PASS & -- & -- & AS4-3 (PASS), \(2.776\times10^{-17}\) & AS4-3
(PASS), -- & symbolic: (d/dx + Delta) sech\^{}2(x/ell) = 0 and (-d\^{}2
+ Delta\^{}2 - Delta') sech tanh = (3/16) sech tanh {[}residuals 0,
0{]} \\
WP2-0a & PASS & \(0\) & \(1.000\times10^{-12}\) & P2-0e (PASS), \(0\) &
YP2-0a (PASS), \(0\) & pair-channel energies of 7 operators in 6
channels (exact algebra of one point) = the table of the Rust program \\
WP2-0b & PASS & \(0\) & \(1.000\times10^{-12}\) & P2-0d (PASS), \(0\) &
YP2-0b (PASS), \(0\) & the three W together = 1/4 of B in every
channel \\
WP2-1 & PASS & \(2.988\times10^{-12}\) & \(1.000\times10^{-8}\) & P2-1a
(PASS), \(6.262\times10^{-14}\) & YP2-1 (PASS), \(7.233\times10^{-13}\)
& self-energies = derivatives of the energy, V\^{}(2), five methods,
random states \\
WP2-2 & PASS & \(2.918\times10^{-11}\) & \(3.000\times10^{-10}\) & P2-2a
(PASS), \(2.919\times10^{-11}\) & YP2-2 (PASS), \(2.918\times10^{-11}\)
& pairing-only limit = AS3 of Part I: /d\^{}3/ = 0.20949330327 (237
iterations) \\
WP2-6b & PASS & \(3.304\times10^{-13}\) & \(1.000\times10^{-10}\) &
P2-6b (PASS), \(1.776\times10^{-15}\) & YP2-6b (PASS),
\(3.268\times10^{-13}\) & the critical pair coupling pi\^{}2/I1(m) =
21.286609756 \\
WP2-6a & PASS & \(6.661\times10^{-15}\) & \(1.000\times10^{-10}\) &
P2-6a (PASS), \(8.882\times10^{-15}\) & YP2-6a (PASS),
\(1.332\times10^{-15}\) & V\^{}(2) from the conditions on the CJT
effective couplings: (G\^{}S, G) = (2.8, 25.2) \\
WP2-6c & PASS & \(7.994\times10^{-15}\) & \(1.000\times10^{-11}\) &
P2-6b (PASS), \(1.776\times10^{-15}\); P2-6c (PASS),
\(3.486\times10^{-14}\) & YP2-6c (PASS), \(1.776\times10^{-15}\) &
effective couplings g\_S, g\_P of 3 potentials x 5 methods = Rust \\
WP2-6d & PASS & \(1.370\times10^{-16}\) & \(1.000\times10^{-11}\) &
P2-6d (PASS), \(0\) & YP2-6d (PASS), \(7.541\times10^{-17}\) &
ground-state energies of 3 potentials x 5 methods = Rust \\
WP2-6e & PASS & \(1.682\times10^{-12}\) & \(1.000\times10^{-8}\) & P2-6d
(PASS), \(0\) & YP2-6e (PASS), \(1.682\times10^{-12}\) & pair amplitudes
/D/ of the ground states = Rust \\
WP2-6f & PASS & \(3.673\times10^{-12}\) & \(1.000\times10^{-8}\) & P2-6d
(PASS), \(0\) & YP2-6f (PASS), \(4.952\times10^{-13}\) & Dirac masses of
the ground states = Rust \\
WP2-6g & PASS & \(3.673\times10^{-12}\) & \(1.000\times10^{-7}\) & P2-6d
(PASS), \(0\) & YP2-6g (PASS), \(4.959\times10^{-13}\) & gaps of the
ground states = Rust \\
WP2-7a HS\_1\_1 & PASS & \(1.603\times10^{-13}\) &
\(1.000\times10^{-9}\) & P2-7a HS\_1\_1 (PASS), \(1.816\times10^{-13}\)
& YP2-7a HS\_1\_1 (PASS), \(1.621\times10^{-13}\) & HS\_1\_1: three
Goldstone modes (the eigenvalues of M(0) equal to 1); the spectrum of
every node is +-E (spread 1.998e-15) \\
WP2-7b HS\_1\_1 & PASS & \(6.573\times10^{-14}\) &
\(1.000\times10^{-7}\) & P2-7a HS\_1\_1 (PASS), \(1.816\times10^{-13}\)
& YP2-7b HS\_1\_1 (PASS), \(5.262\times10^{-14}\) & HS\_1\_1: the
collective modes below the threshold \{0.806128281, 0.813693092\} = Rust
\{0.806128281, 0.813693092\} \\
WP2-7a KS & PASS & \(2.300\times10^{-13}\) & \(1.000\times10^{-9}\) &
P2-7a KS (PASS), \(2.502\times10^{-13}\) & YP2-7a KS (PASS),
\(2.320\times10^{-13}\) & KS: three Goldstone modes (the eigenvalues of
M(0) equal to 1); the spectrum of every node is +-E (spread
1.776e-15) \\
WP2-7b KS & PASS & \(1.482\times10^{-13}\) & \(1.000\times10^{-7}\) &
P2-7a KS (PASS), \(2.502\times10^{-13}\) & YP2-7b KS (PASS),
\(1.333\times10^{-13}\) & KS: the collective modes below the threshold
\{0.603475967, 0.657284727\} = Rust \{0.603475967, 0.657284727\} \\
WP2-7a CJT & PASS & \(2.338\times10^{-13}\) & \(1.000\times10^{-9}\) &
P2-7a CJT (PASS), \(2.485\times10^{-13}\) & YP2-7a CJT (PASS),
\(2.329\times10^{-13}\) & CJT: three Goldstone modes (the eigenvalues of
M(0) equal to 1); the spectrum of every node is +-E (spread
2.442e-15) \\
WP2-7b CJT & PASS & \(1.361\times10^{-13}\) & \(1.000\times10^{-7}\) &
P2-7a CJT (PASS), \(2.485\times10^{-13}\) & YP2-7b CJT (PASS),
\(1.180\times10^{-13}\) & CJT: the collective modes below the threshold
\{0.603475967, 0.644007155\} = Rust \{0.603475967, 0.644007155\} \\
WP2-8a HS\_1\_1 & PASS & \(8.289\times10^{-14}\) &
\(1.000\times10^{-8}\) & P2-8a HS\_1\_1 (PASS), \(2.278\times10^{-15}\);
P2-8b HS\_1\_1 (PASS), \(8.820\times10^{-9}\) & YP2-8a HS\_1\_1 (PASS),
\(3.744\times10^{-15}\) & HS\_1\_1: chi\_11 = 0.005118624664, chi\_FF =
0.02014453922 = Rust; chi\_33 = 8.759e-16 (massless zeugon), no
mixing \\
WP2-8a KS & PASS & \(5.794\times10^{-14}\) & \(1.000\times10^{-8}\) &
P2-8a KS (PASS), \(2.159\times10^{-15}\); P2-8b KS (PASS),
\(4.283\times10^{-9}\) & YP2-8a KS (PASS), \(5.405\times10^{-15}\) & KS:
chi\_11 = 0.004317006163, chi\_FF = 0.02036224303 = Rust; chi\_33 =
4.201e-16 (massless zeugon), no mixing \\
WP2-8a CJT & PASS & \(3.145\times10^{-14}\) & \(1.000\times10^{-8}\) &
P2-8a CJT (PASS), \(1.879\times10^{-15}\); P2-8b CJT (PASS),
\(3.584\times10^{-9}\) & YP2-8a CJT (PASS), \(3.879\times10^{-15}\) &
CJT: chi\_11 = 0.004577873934, chi\_FF = 0.02036224303 = Rust; chi\_33 =
3.067e-16 (massless zeugon), no mixing \\
WP2-8b & PASS & \(9.437\times10^{-16}\) & \(1.000\times10^{-12}\) &
P2-3a (PASS), \(1.475\times10^{-17}\); P2-3b (PASS),
\(5.309\times10^{-15}\) & YP2-8b (PASS), \(1.443\times10^{-14}\) & the
angular averages D0, D1 of the Yukawa kernel in closed form = numerical
integration over the angle \\
WP2-8c & PASS & \(5.919\times10^{-12}\) & \(1.000\times10^{-9}\) & P2-8c
(PASS), \(5.102\times10^{-12}\); P2-8d (PASS), \(2.041\times10^{-11}\) &
YP2-8c (PASS), \(5.918\times10^{-12}\) & the Thouless eigenvalues of the
exchange per unit coupling squared: zeugon W\^{}3 0.03573178317, B
-0.1429271327 = Rust \\
WP2-8d & PASS & \(1.492\times10^{-8}\) & \(1.000\times10^{-7}\) & P2-8f
(PASS), \(0\) & YP2-8d (PASS), \(1.492\times10^{-8}\) & CJT with the
exchange at g\_v = 0.7, g\_F = 0.45: energy -1.324e-4 = Rust (relative;
294 iterations; exchange energy 9.647e-5) \\
WP3-11 HS\_1\_1 & PASS & \(2.492\times10^{-11}\) &
\(1.000\times10^{-7}\) & P3-11 HS\_1\_1 (PASS), \(1.103\times10^{-14}\)
& YP3-11 HS\_1\_1 (PASS), \(9.618\times10^{-12}\) & HS (1, 1): the
accuracy of the Landau form recomputed by LeastSquares (the
coefficients, the minimum, depth and curvature of the polynomial, the
residuals of the fits with Phi\^{}2 .. Phi\^{}(2K), K = 2 .. 6, the fits
on three intervals, the sign changes) = Rust (largest relative
difference) \\
WP3-11 HS\_0\_1 & PASS & \(4.409\times10^{-11}\) &
\(1.000\times10^{-7}\) & P3-11 HS\_0\_1 (PASS), \(5.616\times10^{-13}\)
& YP3-11 HS\_0\_1 (PASS), \(1.478\times10^{-11}\) & HS (0, 1): the
accuracy of the Landau form recomputed by LeastSquares (the
coefficients, the minimum, depth and curvature of the polynomial, the
residuals of the fits with Phi\^{}2 .. Phi\^{}(2K), K = 2 .. 6, the fits
on three intervals, the sign changes) = Rust (largest relative
difference) \\
WP3-11 HS\_0\_0 & PASS & \(4.051\times10^{-11}\) &
\(1.000\times10^{-7}\) & P3-11 HS\_0\_0 (PASS), \(8.612\times10^{-13}\)
& YP3-11 HS\_0\_0 (PASS), \(5.701\times10^{-11}\) & HS (0, 0): the
accuracy of the Landau form recomputed by LeastSquares (the
coefficients, the minimum, depth and curvature of the polynomial, the
residuals of the fits with Phi\^{}2 .. Phi\^{}(2K), K = 2 .. 6, the fits
on three intervals, the sign changes) = Rust (largest relative
difference) \\
WP3-11 KS & PASS & \(3.482\times10^{-9}\) & \(1.000\times10^{-7}\) &
P3-11 KS (PASS), \(7.503\times10^{-13}\) & YP3-11 KS (PASS),
\(3.104\times10^{-9}\) & KS: the accuracy of the Landau form recomputed
by LeastSquares (the coefficients, the minimum, depth and curvature of
the polynomial, the residuals of the fits with Phi\^{}2 .. Phi\^{}(2K),
K = 2 .. 6, the fits on three intervals, the sign changes) = Rust
(largest relative difference) \\
WP3-11 CJT & PASS & \(6.816\times10^{-9}\) & \(1.000\times10^{-7}\) &
P3-11 CJT (PASS), \(2.943\times10^{-13}\) & YP3-11 CJT (PASS),
\(1.216\times10^{-9}\) & CJT: the accuracy of the Landau form recomputed
by LeastSquares (the coefficients, the minimum, depth and curvature of
the polynomial, the residuals of the fits with Phi\^{}2 .. Phi\^{}(2K),
K = 2 .. 6, the fits on three intervals, the sign changes) = Rust
(largest relative difference) \\
WP3-8 & PASS & \(1.554\times10^{-15}\) & \(1.000\times10^{-13}\) & P3-8
(PASS), \(2.442\times10^{-15}\) & YP3-8 (PASS), \(1.332\times10^{-15}\)
& the pair field enters the quasi-particle energies like a mass: the BdG
matrix with N = (M - m) beta (x) 1\_2 and P = Phi\_a (K\^{}3)\^{}dag has
the eigenvalues -+(q\^{}2 + M\^{}2 + Phi\_a\^{}2)\^{}(1/2), eight times
each (20 random q, M, Phi\_a) \\
WP3-10 HS\_1\_1 & PASS & \(2.646\times10^{-7}\) & \(1.000\times10^{-6}\)
& P3-10 HS\_1\_1 (PASS), \(8.786\times10^{-9}\) & YP3-10 HS\_1\_1
(PASS), \(3.072\times10^{-11}\) & HS (1, 1): the coefficient of Phi\^{}2
of the Taylor series at the origin from the first six points of the
curve = the Rust value a\^{}Tay = -0.00616243 (relative) \\
WP3-10 CJT & PASS & \(2.377\times10^{-8}\) & \(1.000\times10^{-6}\) &
P3-10 CJT (PASS), \(7.425\times10^{-8}\) & YP3-10 CJT (PASS),
\(2.872\times10^{-12}\) & CJT: the coefficient of Phi\^{}2 of the Taylor
series at the origin from the first six points of the curve = the Rust
value a\^{}Tay = -0.00429666 (relative) \\
WP3-T1 HS\_1\_1 & PASS & \(7.903\times10^{-14}\) &
\(1.000\times10^{-11}\) & P3-T1 HS\_1\_1 (PASS), \(2.357\times10^{-14}\)
& YP3-T1 HS\_1\_1 (PASS), \(6.164\times10^{-14}\) & HS (1, 1): the
thermal states at T = 0.005 reproduce the ground state of T = 0 (/change
of Phi0/ + relative change of the energy) \\
WP3-T1 CJT & PASS & \(9.391\times10^{-14}\) & \(1.000\times10^{-11}\) &
P3-T1 CJT (PASS), \(8.078\times10^{-14}\) & YP3-T1 CJT (PASS),
\(1.046\times10^{-13}\) & CJT: the thermal states at T = 0.005 reproduce
the ground state of T = 0 (/change of Phi0/ + relative change of the
energy) \\
WP3-T3 HS\_1\_1 & PASS & \(7.291\times10^{-14}\) &
\(1.000\times10^{-10}\) & P3-T3 HS\_1\_1 (PASS), \(4.408\times10^{-14}\)
& YP3-T3 HS\_1\_1 (PASS), \(6.756\times10^{-14}\); YP3-T9 HS\_1\_1
(PASS), \(2.420\times10^{-14}\) & HS (1, 1): the critical temperature
T\_c = 0.2604672394 (bisection of lambda\^{}P(T) = 1) = Rust (relative),
and lambda\^{}P(0) = 1.183842814 = Rust; the Dirac mass of the normal
state of the thermal engine at T\_c = 0.2066851932 = the scalar
self-consistency \\
WP3-T3 HS\_0\_1 & PASS & \(4.653\times10^{-13}\) &
\(1.000\times10^{-10}\) & P3-T3 HS\_0\_1 (PASS), \(4.028\times10^{-13}\)
& YP3-T9 HS\_0\_1 (PASS), \(3.464\times10^{-14}\) & HS (0, 1): the
critical temperature T\_c = 0.2404654563 (bisection of lambda\^{}P(T) =
1) = Rust (relative), and lambda\^{}P(0) = 1.150958292 = Rust \\
WP3-T3 HS\_0\_0 & PASS & \(2.056\times10^{-13}\) &
\(1.000\times10^{-10}\) & P3-T3 HS\_0\_0 (PASS), \(1.976\times10^{-13}\)
& YP3-T9 HS\_0\_0 (PASS), \(1.688\times10^{-14}\) & HS (0, 0): the
critical temperature T\_c = 0.2137990984 (bisection of lambda\^{}P(T) =
1) = Rust (relative), and lambda\^{}P(0) = 1.11337598 = Rust \\
WP3-T3 KS & PASS & \(3.755\times10^{-13}\) & \(1.000\times10^{-10}\) &
P3-T3 KS (PASS), \(1.987\times10^{-13}\) & YP3-T9 KS (PASS),
\(5.906\times10^{-14}\) & KS: the critical temperature T\_c =
0.2221187371 (bisection of lambda\^{}P(T) = 1) = Rust (relative), and
lambda\^{}P(0) = 1.11337598 = Rust; the Dirac mass of the normal state
of the thermal engine at T\_c = 0.1622929325 = the scalar
self-consistency \\
WP3-T3 CJT & PASS & \(3.755\times10^{-13}\) & \(1.000\times10^{-10}\) &
P3-T3 CJT (PASS), \(1.978\times10^{-13}\) & YP3-T3 CJT (PASS),
\(2.138\times10^{-13}\); YP3-T9 CJT (PASS), \(5.906\times10^{-14}\) &
CJT: the critical temperature T\_c = 0.2221187371 (bisection of
lambda\^{}P(T) = 1) = Rust (relative), and lambda\^{}P(0) = 1.11337598 =
Rust; the Dirac mass of the normal state of the thermal engine at T\_c =
0.1622929325 = the scalar self-consistency \\
WP3-T9 HS\_1\_1 & PASS & \(0\) & \(1.000\times10^{-12}\) & P3-T9
HS\_1\_1 (PASS), \(0\) & YP3-T9 HS\_1\_1 (PASS), \(2.420\times10^{-14}\)
& HS (1, 1): lambda\^{}P(T) of the normal state does not increase on 60
temperatures from 0.01 to 0.6 (from 1.18384 to 0.620857; the largest
rise), so T\_c is the only solution of lambda\^{}P(T) = 1 there \\
WP3-T9 HS\_0\_1 & PASS & \(0\) & \(1.000\times10^{-12}\) & P3-T9
HS\_0\_1 (PASS), \(0\) & YP3-T9 HS\_0\_1 (PASS), \(3.464\times10^{-14}\)
& HS (0, 1): lambda\^{}P(T) of the normal state does not increase on 60
temperatures from 0.01 to 0.6 (from 1.15096 to 0.60469; the largest
rise), so T\_c is the only solution of lambda\^{}P(T) = 1 there \\
WP3-T9 HS\_0\_0 & PASS & \(0\) & \(1.000\times10^{-12}\) & P3-T9
HS\_0\_0 (PASS), \(0\) & YP3-T9 HS\_0\_0 (PASS), \(1.688\times10^{-14}\)
& HS (0, 0): lambda\^{}P(T) of the normal state does not increase on 60
temperatures from 0.01 to 0.6 (from 1.11338 to 0.584946; the largest
rise), so T\_c is the only solution of lambda\^{}P(T) = 1 there \\
WP3-T9 KS & PASS & \(0\) & \(1.000\times10^{-12}\) & P3-T9 KS (PASS),
\(0\) & YP3-T9 KS (PASS), \(5.906\times10^{-14}\) & KS: lambda\^{}P(T)
of the normal state does not increase on 60 temperatures from 0.01 to
0.6 (from 1.11338 to 0.588434; the largest rise), so T\_c is the only
solution of lambda\^{}P(T) = 1 there \\
WP3-T9 CJT & PASS & \(0\) & \(1.000\times10^{-12}\) & P3-T9 CJT (PASS),
\(0\) & YP3-T9 CJT (PASS), \(5.906\times10^{-14}\) & CJT: lambda\^{}P(T)
of the normal state does not increase on 60 temperatures from 0.01 to
0.6 (from 1.11338 to 0.588434; the largest rise), so T\_c is the only
solution of lambda\^{}P(T) = 1 there \\
WP3-T2 CJT & PASS & \(1.589\times10^{-12}\) & \(1.000\times10^{-8}\) &
P3-T2 CJT (PASS), \(1.142\times10^{-12}\) & YP3-T2 CJT (PASS),
\(6.209\times10^{-13}\) & CJT: the superfluid at 0.5 T\_c and 0.9 T\_c
(thermal engine): the order parameter and the free energy = Rust
(largest difference; the free energy relative), and the gap equation
holds (largest /g\^{}P I1/pi\^{}2 - 1/; the sum of both) \\
WP3-T8 CJT & PASS & \(2.765\times10^{-13}\) & \(1.000\times10^{-10}\) &
P3-T8 CJT (PASS), \(3.271\times10^{-13}\) & YP3-T8 CJT (PASS),
\(1.717\times10^{-13}\) & CJT: the Dirac mass of the superfluid at 0.5
T\_c and 0.9 T\_c equals that at T = 0, 0.1622929325, and the Rust value
(largest difference) \\
WP3-T2 HS\_0\_0 & PASS & \(9.924\times10^{-12}\) &
\(1.000\times10^{-9}\) & P3-T2 HS\_0\_0 (PASS), \(1.214\times10^{-12}\)
& YP3-T4 (PASS), \(9.807\times10^{-12}\) & HS (0, 0): the order
parameter of every temperature of the table from the gap equation alone
(with the Dirac mass m\_LL, without the self-consistent engine) = Rust
(largest difference) \\
WP3-T5 CJT & PASS & \(1.431\times10^{-7}\) & \(1.000\times10^{-5}\) &
P3-T5 CJT (PASS), \(2.348\times10^{-1}\) & YP3-T5 CJT (PASS),
\(3.080\times10^{-7}\) & CJT at 0.9 T\_c: the Taylor coefficients a =
-0.0010510272, b = 0.014755098 of the free energy at the origin and the
error 2.277e-2 of the order parameter (-a/2b)\^{}(1/2) of Landau's
theory (0.228 times 1 - T/T\_c) = Rust (largest relative difference of a
and b, difference of the error) \\
WP3-T6 HS\_1\_1 & PASS & \(5.525\times10^{-11}\) &
\(1.000\times10^{-7}\) & P3-T6 HS\_1\_1 (PASS), \(0\) & YP3-T6 HS\_1\_1
(PASS), \(7.060\times10^{-11}\) & HS (1, 1): the largest residuals of
the Landau form and of the quartic form at T/T\_c = 0, 0.5, 0.8, 0.9,
0.95, 0.98, recomputed by LeastSquares from the curves of the Rust
program = Rust (largest relative difference); the quartic residual falls
towards T\_c (1.828e-2, 8.811e-3, 3.734e-3 of the depth at 0.9, 0.95,
0.98 T\_c) and that of the Landau form is 1.668e-5 at 0.98 T\_c: yes \\
WP3-T7 HS\_1\_1 & PASS & \(5.804\times10^{-12}\) &
\(1.000\times10^{-8}\) & P3-T7 HS\_1\_1 (PASS), \(5.579\times10^{-12}\)
& YP3-T7 HS\_1\_1 (PASS), \(5.638\times10^{-12}\) & HS (1, 1): at T = 0
the free energy is the effective potential: the residual of the Landau
form of the curve at T = 0 is that of Section III.L (relative
difference) \\
WP3-T6 HS\_0\_1 & PASS & \(2.410\times10^{-11}\) &
\(1.000\times10^{-7}\) & P3-T6 HS\_0\_1 (PASS), \(0\) & YP3-T6 HS\_0\_1
(PASS), \(9.270\times10^{-12}\) & HS (0, 1): the largest residuals of
the Landau form and of the quartic form at T/T\_c = 0, 0.5, 0.8, 0.9,
0.95, 0.98, recomputed by LeastSquares from the curves of the Rust
program = Rust (largest relative difference); the quartic residual falls
towards T\_c (1.597e-2, 7.668e-3, 3.244e-3 of the depth at 0.9, 0.95,
0.98 T\_c) and that of the Landau form is 1.252e-5 at 0.98 T\_c: yes \\
WP3-T7 HS\_0\_1 & PASS & \(2.383\times10^{-12}\) &
\(1.000\times10^{-8}\) & P3-T7 HS\_0\_1 (PASS), \(5.897\times10^{-13}\)
& YP3-T7 HS\_0\_1 (PASS), \(1.304\times10^{-12}\) & HS (0, 1): at T = 0
the free energy is the effective potential: the residual of the Landau
form of the curve at T = 0 is that of Section III.L (relative
difference) \\
WP3-T6 HS\_0\_0 & PASS & \(1.579\times10^{-10}\) &
\(1.000\times10^{-7}\) & P3-T6 HS\_0\_0 (PASS), \(0\) & YP3-T6 HS\_0\_0
(PASS), \(4.884\times10^{-10}\) & HS (0, 0): the largest residuals of
the Landau form and of the quartic form at T/T\_c = 0, 0.5, 0.8, 0.9,
0.95, 0.98, recomputed by LeastSquares from the curves of the Rust
program = Rust (largest relative difference); the quartic residual falls
towards T\_c (1.267e-2, 6.063e-3, 2.562e-3 of the depth at 0.9, 0.95,
0.98 T\_c) and that of the Landau form is 7.750e-6 at 0.98 T\_c: yes \\
WP3-T7 HS\_0\_0 & PASS & \(2.475\times10^{-11}\) &
\(1.000\times10^{-8}\) & P3-T7 HS\_0\_0 (PASS), \(1.841\times10^{-11}\)
& YP3-T7 HS\_0\_0 (PASS), \(1.891\times10^{-11}\) & HS (0, 0): at T = 0
the free energy is the effective potential: the residual of the Landau
form of the curve at T = 0 is that of Section III.L (relative
difference) \\
WP3-T6 KS & PASS & \(5.965\times10^{-11}\) & \(1.000\times10^{-7}\) &
P3-T6 KS (PASS), \(0\) & YP3-T6 KS (PASS), \(1.306\times10^{-10}\) & KS:
the largest residuals of the Landau form and of the quartic form at
T/T\_c = 0, 0.5, 0.8, 0.9, 0.95, 0.98, recomputed by LeastSquares from
the curves of the Rust program = Rust (largest relative difference); the
quartic residual falls towards T\_c (1.579e-2, 7.470e-3, 3.134e-3 of the
depth at 0.9, 0.95, 0.98 T\_c) and that of the Landau form is 9.969e-6
at 0.98 T\_c: yes \\
WP3-T7 KS & PASS & \(4.070\times10^{-11}\) & \(1.000\times10^{-8}\) &
P3-T7 KS (PASS), \(3.458\times10^{-11}\) & YP3-T7 KS (PASS),
\(3.539\times10^{-11}\) & KS: at T = 0 the free energy is the effective
potential: the residual of the Landau form of the curve at T = 0 is that
of Section III.L (relative difference) \\
WP3-T6 CJT & PASS & \(8.948\times10^{-12}\) & \(1.000\times10^{-7}\) &
P3-T6 CJT (PASS), \(0\) & YP3-T6 CJT (PASS), \(5.804\times10^{-11}\) &
CJT: the largest residuals of the Landau form and of the quartic form at
T/T\_c = 0, 0.5, 0.8, 0.9, 0.95, 0.98, recomputed by LeastSquares from
the curves of the Rust program = Rust (largest relative difference); the
quartic residual falls towards T\_c (1.579e-2, 7.470e-3, 3.134e-3 of the
depth at 0.9, 0.95, 0.98 T\_c) and that of the Landau form is 9.969e-6
at 0.98 T\_c: yes \\
WP3-T7 CJT & PASS & \(4.876\times10^{-11}\) & \(1.000\times10^{-8}\) &
P3-T7 CJT (PASS), \(4.596\times10^{-11}\) & YP3-T7 CJT (PASS),
\(4.636\times10^{-11}\) & CJT: at T = 0 the free energy is the effective
potential: the residual of the Landau form of the curve at T = 0 is that
of Section III.L (relative difference) \\
WP3-G1 HS\_1\_1 & PASS & \(7.975\times10^{-10}\) &
\(1.000\times10^{-7}\) & P3-G1 HS\_1\_1 (PASS), \(7.975\times10^{-10}\)
& YP3-G1 HS\_1\_1 (PASS), \(7.976\times10^{-10}\) & HS (1, 1): at 0.9
T\_c the coefficient a(0) of /Phi/\^{}2 from the static pair response
(-1.343e-3) equals the Taylor coefficient a\^{}Tay of the nonlinear free
energy of Section III.T (-1.343e-3) (relative difference) \\
WP3-G2 HS\_1\_1 & PASS & \(1.099\times10^{-12}\) &
\(1.000\times10^{-9}\) & P3-G2 HS\_1\_1 (PASS), \(6.276\times10^{-13}\)
& YP3-G2 HS\_1\_1 (PASS), \(9.724\times10^{-12}\) & HS (1, 1): a(0, T) =
(1 - lambda\^{}P(T))/(lambda\^{}P(T) g\^{}P) at T/T\_c = 0.9, 0.998,
1.002, with lambda\^{}P(T) of Section III.T.4 (largest relative
difference) \\
WP3-G10 HS\_1\_1 & PASS & \(1.496\times10^{-12}\) &
\(1.000\times10^{-8}\) & P3-G10 HS\_1\_1 (PASS), \(5.463\times10^{-13}\)
& YP3-G10 HS\_1\_1 (PASS), \(4.438\times10^{-13}\) & HS (1, 1): alpha =
T\_c da/dT = 0.01403788523 from central differences of the static pair
response (Richardson) = 0.01403788523 in closed form (lambda\^{}P(T\_c)
= 1., M\_n = 0.2066851932) = Rust (largest relative difference) \\
WP3-G3 HS\_1\_1 & PASS & \(4.080\times10^{-14}\) &
\(1.000\times10^{-9}\) & P3-G3 HS\_1\_1 (PASS), \(2.197\times10^{-15}\)
& YP3-G3 HS\_1\_1 (PASS), \(7.218\times10^{-14}\) & HS (1, 1): a(0,
T\_c) vanishes to the accuracy of T\_c (/a(0, T\_c)/ / alpha) \\
WP3-G4 HS\_1\_1 & PASS & \(0\) & \(1.000\times10^{-12}\) & P3-G4
HS\_1\_1 (PASS), \(1.089\times10^{-33}\) & YP3-G4 HS\_1\_1 (PASS), \(0\)
& HS (1, 1): at zero pair momentum no other pair channel is induced
(norm of the self-consistent components of P\^{}a other than
(K\^{}iota)\^{}dag) \\
WP3-G1 HS\_0\_1 & PASS & \(1.747\times10^{-10}\) &
\(1.000\times10^{-7}\) & P3-G1 HS\_0\_1 (PASS), \(1.747\times10^{-10}\)
& YP3-G1 HS\_0\_1 (PASS), \(1.748\times10^{-10}\) & HS (0, 1): at 0.9
T\_c the coefficient a(0) of /Phi/\^{}2 from the static pair response
(-1.209e-3) equals the Taylor coefficient a\^{}Tay of the nonlinear free
energy of Section III.T (-1.209e-3) (relative difference) \\
WP3-G2 HS\_0\_1 & PASS & \(7.142\times10^{-13}\) &
\(1.000\times10^{-9}\) & P3-G2 HS\_0\_1 (PASS), \(3.926\times10^{-12}\)
& YP3-G2 HS\_0\_1 (PASS), \(7.318\times10^{-12}\) & HS (0, 1): a(0, T) =
(1 - lambda\^{}P(T))/(lambda\^{}P(T) g\^{}P) at T/T\_c = 0.9, 0.998,
1.002, with lambda\^{}P(T) of Section III.T.4 (largest relative
difference) \\
WP3-G10 HS\_0\_1 & PASS & \(2.014\times10^{-12}\) &
\(1.000\times10^{-8}\) & P3-G10 HS\_0\_1 (PASS), \(2.118\times10^{-12}\)
& YP3-G10 HS\_0\_1 (PASS), \(1.914\times10^{-12}\) & HS (0, 1): alpha =
T\_c da/dT = 0.01269913755 from central differences of the static pair
response (Richardson) = 0.01269913755 in closed form (lambda\^{}P(T\_c)
= 1., M\_n = 0.2) = Rust (largest relative difference) \\
WP3-G3 HS\_0\_1 & PASS & \(4.639\times10^{-14}\) &
\(1.000\times10^{-9}\) & P3-G3 HS\_0\_1 (PASS), \(9.010\times10^{-15}\)
& YP3-G3 HS\_0\_1 (PASS), \(1.071\times10^{-13}\) & HS (0, 1): a(0,
T\_c) vanishes to the accuracy of T\_c (/a(0, T\_c)/ / alpha) \\
WP3-G4 HS\_0\_1 & PASS & \(4.441\times10^{-16}\) &
\(1.000\times10^{-12}\) & P3-G4 HS\_0\_1 (PASS), \(1.929\times10^{-16}\)
& YP3-G4 HS\_0\_1 (PASS), \(1.392\times10^{-18}\) & HS (0, 1): at zero
pair momentum no other pair channel is induced (norm of the
self-consistent components of P\^{}a other than (K\^{}iota)\^{}dag) \\
WP3-G1 HS\_0\_0 & PASS & \(8.823\times10^{-10}\) &
\(1.000\times10^{-7}\) & P3-G1 HS\_0\_0 (PASS), \(8.823\times10^{-10}\)
& YP3-G1 HS\_0\_0 (PASS), \(8.823\times10^{-10}\) & HS (0, 0): at 0.9
T\_c the coefficient a(0) of /Phi/\^{}2 from the static pair response
(-1.016e-3) equals the Taylor coefficient a\^{}Tay of the nonlinear free
energy of Section III.T (-1.016e-3) (relative difference) \\
WP3-G2 HS\_0\_0 & PASS & \(4.362\times10^{-13}\) &
\(1.000\times10^{-9}\) & P3-G2 HS\_0\_0 (PASS), \(5.015\times10^{-12}\)
& YP3-G2 HS\_0\_0 (PASS), \(1.047\times10^{-11}\) & HS (0, 0): a(0, T) =
(1 - lambda\^{}P(T))/(lambda\^{}P(T) g\^{}P) at T/T\_c = 0.9, 0.998,
1.002, with lambda\^{}P(T) of Section III.T.4 (largest relative
difference) \\
WP3-G10 HS\_0\_0 & PASS & \(2.713\times10^{-12}\) &
\(1.000\times10^{-8}\) & P3-G10 HS\_0\_0 (PASS), \(2.378\times10^{-12}\)
& YP3-G10 HS\_0\_0 (PASS), \(2.272\times10^{-12}\) & HS (0, 0): alpha =
T\_c da/dT = 0.01074338535 from central differences of the static pair
response (Richardson) = 0.01074338535 in closed form (lambda\^{}P(T\_c)
= 1., M\_n = 0.2) = Rust (largest relative difference) \\
WP3-G3 HS\_0\_0 & PASS & \(5.930\times10^{-14}\) &
\(1.000\times10^{-9}\) & P3-G3 HS\_0\_0 (PASS), \(2.071\times10^{-15}\)
& YP3-G3 HS\_0\_0 (PASS), \(5.145\times10^{-14}\) & HS (0, 0): a(0,
T\_c) vanishes to the accuracy of T\_c (/a(0, T\_c)/ / alpha) \\
WP3-G4 HS\_0\_0 & PASS & \(3.259\times10^{-18}\) &
\(1.000\times10^{-12}\) & P3-G4 HS\_0\_0 (PASS), \(1.116\times10^{-16}\)
& YP3-G4 HS\_0\_0 (PASS), \(1.063\times10^{-18}\) & HS (0, 0): at zero
pair momentum no other pair channel is induced (norm of the
self-consistent components of P\^{}a other than (K\^{}iota)\^{}dag) \\
WP3-G1 KS & PASS & \(4.374\times10^{-10}\) & \(1.000\times10^{-7}\) &
P3-G1 KS (PASS), \(4.374\times10^{-10}\) & YP3-G1 KS (PASS),
\(4.375\times10^{-10}\) & KS: at 0.9 T\_c the coefficient a(0) of
/Phi/\^{}2 from the static pair response (-1.051e-3) equals the Taylor
coefficient a\^{}Tay of the nonlinear free energy of Section III.T
(-1.051e-3) (relative difference) \\
WP3-G2 KS & PASS & \(5.240\times10^{-12}\) & \(1.000\times10^{-9}\) &
P3-G2 KS (PASS), \(1.572\times10^{-12}\) & YP3-G2 KS (PASS),
\(1.635\times10^{-11}\) & KS: a(0, T) = (1 -
lambda\^{}P(T))/(lambda\^{}P(T) g\^{}P) at T/T\_c = 0.9, 0.998, 1.002,
with lambda\^{}P(T) of Section III.T.4 (largest relative difference) \\
WP3-G10 KS & PASS & \(6.309\times10^{-12}\) & \(1.000\times10^{-8}\) &
P3-G10 KS (PASS), \(4.843\times10^{-12}\) & YP3-G10 KS (PASS),
\(4.420\times10^{-12}\) & KS: alpha = T\_c da/dT = 0.01117552279 from
central differences of the static pair response (Richardson) =
0.01117552279 in closed form (lambda\^{}P(T\_c) = 1., M\_n =
0.1622929325) = Rust (largest relative difference) \\
WP3-G3 KS & PASS & \(4.443\times10^{-14}\) & \(1.000\times10^{-9}\) &
P3-G3 KS (PASS), \(2.934\times10^{-15}\) & YP3-G3 KS (PASS),
\(2.192\times10^{-13}\) & KS: a(0, T\_c) vanishes to the accuracy of
T\_c (/a(0, T\_c)/ / alpha) \\
WP3-G4 KS & PASS & \(4.441\times10^{-16}\) & \(1.000\times10^{-12}\) &
P3-G4 KS (PASS), \(3.135\times10^{-16}\) & YP3-G4 KS (PASS),
\(3.134\times10^{-16}\) & KS: at zero pair momentum no other pair
channel is induced (norm of the self-consistent components of P\^{}a
other than (K\^{}iota)\^{}dag) \\
WP3-G1 CJT & PASS & \(4.421\times10^{-10}\) & \(1.000\times10^{-7}\) &
P3-G1 CJT (PASS), \(4.421\times10^{-10}\) & YP3-G1 CJT (PASS),
\(4.422\times10^{-10}\) & CJT: at 0.9 T\_c the coefficient a(0) of
/Phi/\^{}2 from the static pair response (-1.051e-3) equals the Taylor
coefficient a\^{}Tay of the nonlinear free energy of Section III.T
(-1.051e-3) (relative difference) \\
WP3-G2 CJT & PASS & \(5.240\times10^{-12}\) & \(1.000\times10^{-9}\) &
P3-G2 CJT (PASS), \(1.258\times10^{-12}\) & YP3-G2 CJT (PASS),
\(1.583\times10^{-11}\) & CJT: a(0, T) = (1 -
lambda\^{}P(T))/(lambda\^{}P(T) g\^{}P) at T/T\_c = 0.9, 0.998, 1.002,
with lambda\^{}P(T) of Section III.T.4 (largest relative difference) \\
WP3-G10 CJT & PASS & \(6.950\times10^{-12}\) & \(1.000\times10^{-8}\) &
P3-G10 CJT (PASS), \(5.127\times10^{-12}\) & YP3-G10 CJT (PASS),
\(4.702\times10^{-12}\) & CJT: alpha = T\_c da/dT = 0.01117552279 from
central differences of the static pair response (Richardson) =
0.01117552279 in closed form (lambda\^{}P(T\_c) = 1., M\_n =
0.1622929325) = Rust (largest relative difference) \\
WP3-G3 CJT & PASS & \(4.527\times10^{-14}\) & \(1.000\times10^{-9}\) &
P3-G3 CJT (PASS), \(2.725\times10^{-15}\) & YP3-G3 CJT (PASS),
\(2.192\times10^{-13}\) & CJT: a(0, T\_c) vanishes to the accuracy of
T\_c (/a(0, T\_c)/ / alpha) \\
WP3-G4 CJT & PASS & \(4.441\times10^{-16}\) & \(1.000\times10^{-12}\) &
P3-G4 CJT (PASS), \(2.218\times10^{-16}\) & YP3-G4 CJT (PASS),
\(3.134\times10^{-16}\) & CJT: at zero pair momentum no other pair
channel is induced (norm of the self-consistent components of P\^{}a
other than (K\^{}iota)\^{}dag) \\
WP3-G5 HS\_1\_1 & PASS & \(4.978\times10^{-12}\) &
\(1.000\times10^{-8}\) & P3-G5 HS\_1\_1 (PASS), \(2.845\times10^{-11}\)
& YP3-G5 HS\_1\_1 (PASS), \(3.656\times10^{-10}\) & HS (1, 1): the
isospin twist gives the gradient coefficient of the phase twist, c =
0.01548747577 (isospin invariance; relative difference) \\
WP3-G6 HS\_1\_1 & PASS & \(1.064\times10^{-11}\) &
\(1.000\times10^{-8}\) & P3-G6 HS\_1\_1 (PASS), \(0\) & YP3-G6 HS\_1\_1
(PASS), -- & HS (1, 1): c \textgreater{} 0 and the extrapolation in k
converges: c(0.01) = 0.01548761918, c(0.02) = 0.01548804939,
extrapolated 0.01548747577 = Rust (relative difference; 1 is added if c
\textless= 0 or the two values differ by more than 1e-3 of c) \\
WP3-G5 HS\_0\_1 & PASS & \(2.030\times10^{-11}\) &
\(1.000\times10^{-8}\) & P3-G5 HS\_0\_1 (PASS), \(3.195\times10^{-12}\)
& YP3-G5 HS\_0\_1 (PASS), \(2.265\times10^{-10}\) & HS (0, 1): the
isospin twist gives the gradient coefficient of the phase twist, c =
0.01713758647 (isospin invariance; relative difference) \\
WP3-G6 HS\_0\_1 & PASS & \(2.462\times10^{-11}\) &
\(1.000\times10^{-8}\) & P3-G6 HS\_0\_1 (PASS), \(0\) & YP3-G6 HS\_0\_1
(PASS), -- & HS (0, 1): c \textgreater{} 0 and the extrapolation in k
converges: c(0.01) = 0.01713773698, c(0.02) = 0.01713818851,
extrapolated 0.01713758647 = Rust (relative difference; 1 is added if c
\textless= 0 or the two values differ by more than 1e-3 of c) \\
WP3-G5 HS\_0\_0 & PASS & \(1.522\times10^{-11}\) &
\(1.000\times10^{-8}\) & P3-G5 HS\_0\_0 (PASS), \(3.099\times10^{-12}\)
& YP3-G5 HS\_0\_0 (PASS), \(1.692\times10^{-10}\) & HS (0, 0): the
isospin twist gives the gradient coefficient of the phase twist, c =
0.01952798529 (isospin invariance; relative difference) \\
WP3-G6 HS\_0\_0 & PASS & \(3.960\times10^{-11}\) &
\(1.000\times10^{-8}\) & P3-G6 HS\_0\_0 (PASS), \(0\) & YP3-G6 HS\_0\_0
(PASS), -- & HS (0, 0): c \textgreater{} 0 and the extrapolation in k
converges: c(0.01) = 0.01952812997, c(0.02) = 0.019528564, extrapolated
0.01952798529 = Rust (relative difference; 1 is added if c \textless= 0
or the two values differ by more than 1e-3 of c) \\
WP3-G5 KS & PASS & \(9.723\times10^{-12}\) & \(1.000\times10^{-8}\) &
P3-G5 KS (PASS), \(1.707\times10^{-11}\) & YP3-G5 KS (PASS),
\(6.579\times10^{-11}\) & KS: the isospin twist gives the gradient
coefficient of the phase twist, c = 0.01927177616 (isospin invariance;
relative difference) \\
WP3-G6 KS & PASS & \(5.371\times10^{-11}\) & \(1.000\times10^{-8}\) &
P3-G6 KS (PASS), \(0\) & YP3-G6 KS (PASS), -- & KS: c \textgreater{} 0
and the extrapolation in k converges: c(0.01) = 0.0192719291, c(0.02) =
0.01927238792, extrapolated 0.01927177616 = Rust (relative difference; 1
is added if c \textless= 0 or the two values differ by more than 1e-3 of
c) \\
WP3-G5 CJT & PASS & \(1.491\times10^{-11}\) & \(1.000\times10^{-8}\) &
P3-G5 CJT (PASS), \(1.327\times10^{-11}\) & YP3-G5 CJT (PASS),
\(6.741\times10^{-11}\) & CJT: the isospin twist gives the gradient
coefficient of the phase twist, c = 0.01927177616 (isospin invariance;
relative difference) \\
WP3-G6 CJT & PASS & \(4.072\times10^{-11}\) & \(1.000\times10^{-8}\) &
P3-G6 CJT (PASS), \(0\) & YP3-G6 CJT (PASS), -- & CJT: c \textgreater{}
0 and the extrapolation in k converges: c(0.01) = 0.0192719291, c(0.02)
= 0.01927238792, extrapolated 0.01927177616 = Rust (relative difference;
1 is added if c \textless= 0 or the two values differ by more than 1e-3
of c) \\
WP3-G9 HS\_1\_1 & PASS & \(3.527\times10^{-10}\) &
\(1.000\times10^{-7}\) & P3-G9 HS\_1\_1 (PASS), \(0\) & YP3-G9 HS\_1\_1
(PASS), -- & HS (1, 1): with the cutoff on the momenta of the two
charges a(k) - a(0) is linear in /k/ (ratio 2.01391 for k = 0.01 and
0.005), with the cutoff on the relative momentum quadratic (ratio
4.000035); the four differences = Rust (largest relative difference; 1
is added if a ratio is off by more than 0.02 of 2 or 1e-3 of 4) \\
WP3-G9 HS\_0\_1 & PASS & \(1.890\times10^{-10}\) &
\(1.000\times10^{-7}\) & P3-G9 HS\_0\_1 (PASS), \(0\) & YP3-G9 HS\_0\_1
(PASS), -- & HS (0, 1): with the cutoff on the momenta of the two
charges a(k) - a(0) is linear in /k/ (ratio 2.01448 for k = 0.01 and
0.005), with the cutoff on the relative momentum quadratic (ratio
4.000036); the four differences = Rust (largest relative difference; 1
is added if a ratio is off by more than 0.02 of 2 or 1e-3 of 4) \\
WP3-G9 HS\_0\_0 & PASS & \(1.919\times10^{-10}\) &
\(1.000\times10^{-7}\) & P3-G9 HS\_0\_0 (PASS), \(0\) & YP3-G9 HS\_0\_0
(PASS), -- & HS (0, 0): with the cutoff on the momenta of the two
charges a(k) - a(0) is linear in /k/ (ratio 2.01532 for k = 0.01 and
0.005), with the cutoff on the relative momentum quadratic (ratio
4.000035); the four differences = Rust (largest relative difference; 1
is added if a ratio is off by more than 0.02 of 2 or 1e-3 of 4) \\
WP3-G9 KS & PASS & \(2.783\times10^{-10}\) & \(1.000\times10^{-7}\) &
P3-G9 KS (PASS), \(0\) & YP3-G9 KS (PASS), -- & KS: with the cutoff on
the momenta of the two charges a(k) - a(0) is linear in /k/ (ratio
2.0152 for k = 0.01 and 0.005), with the cutoff on the relative momentum
quadratic (ratio 4.000036); the four differences = Rust (largest
relative difference; 1 is added if a ratio is off by more than 0.02 of 2
or 1e-3 of 4) \\
WP3-G9 CJT & PASS & \(1.410\times10^{-10}\) & \(1.000\times10^{-7}\) &
P3-G9 CJT (PASS), \(0\) & YP3-G9 CJT (PASS), -- & CJT: with the cutoff
on the momenta of the two charges a(k) - a(0) is linear in /k/ (ratio
2.0152 for k = 0.01 and 0.005), with the cutoff on the relative momentum
quadratic (ratio 4.000036); the four differences = Rust (largest
relative difference; 1 is added if a ratio is off by more than 0.02 of 2
or 1e-3 of 4) \\
WP3-G8 & PASS & \(1.609\times10^{-10}\) & \(1.000\times10^{-8}\) & P3-G8
(PASS), \(1.874\times10^{-10}\) & YP3-G8 CJT (PASS),
\(2.252\times10^{-10}\) & CJT: the gradient coefficient does not depend
on the quadrature: six directions and 24 radial nodes 0.0192717761556,
144 directions 0.0192717761587, 40 radial nodes 0.0192717761551 (largest
relative difference) \\
WP3-G12 & PASS & \(6.308\times10^{-8}\) & \(1.000\times10^{-6}\) &
P3-G12 CJT (PASS), \(9.245\times10^{-8}\) & YP3-G12 CJT (PASS),
\(1.715\times10^{-7}\) & CJT at T\_c: the quartic coefficient b =
0.01472373584 of the free energy at the origin (four applied fields,
thermal engine) is positive and = Rust 0.01472373676 (relative
difference; a(T\_c) = 5.619e-14) \\
WP3-G11 HS\_1\_1 & PASS & \(3.247\times10^{-11}\) &
\(1.000\times10^{-6}\) & P3-G11 HS\_1\_1 (PASS), \(4.441\times10^{-16}\)
& YP3-G11 HS\_1\_1 (PASS), \(1.274\times10^{-10}\) & HS (1, 1): the
Ginzburg number by the first form 11.379614 and by the second form
11.379614 = Rust 11.379614 (largest relative difference) \\
WP3-G11 HS\_0\_1 & PASS & \(7.599\times10^{-11}\) &
\(1.000\times10^{-6}\) & P3-G11 HS\_0\_1 (PASS), \(2.220\times10^{-16}\)
& YP3-G11 HS\_0\_1 (PASS), \(6.345\times10^{-10}\) & HS (0, 1): the
Ginzburg number by the first form 9.7677449 and by the second form
9.7677449 = Rust 9.7677449 (largest relative difference) \\
WP3-G11 HS\_0\_0 & PASS & \(1.212\times10^{-10}\) &
\(1.000\times10^{-6}\) & P3-G11 HS\_0\_0 (PASS), \(6.661\times10^{-16}\)
& YP3-G11 HS\_0\_0 (PASS), \(2.048\times10^{-10}\) & HS (0, 0): the
Ginzburg number by the first form 8.2771765 and by the second form
8.2771765 = Rust 8.2771765 (largest relative difference) \\
WP3-G11 KS & PASS & \(1.563\times10^{-10}\) & \(1.000\times10^{-6}\) &
P3-G11 KS (PASS), \(0\) & YP3-G11 KS (PASS), \(7.051\times10^{-10}\) &
KS: the Ginzburg number by the first form 7.6207226 and by the second
form 7.6207226 = Rust 7.6207226 (largest relative difference) \\
WP3-G11 CJT & PASS & \(1.260\times10^{-7}\) & \(1.000\times10^{-6}\) &
P3-G11 CJT (PASS), \(4.441\times10^{-16}\) & YP3-G11 CJT (PASS),
\(6.806\times10^{-10}\) & CJT: the Ginzburg number by the first form
7.6207222 and by the second form 7.6207222 = Rust 7.6207232 (largest
relative difference) \\
WP3-H2 HS\_1\_1 & PASS & \(3.153\times10^{-14}\) &
\(1.000\times10^{-10}\) & P3-H2 HS\_1\_1 (PASS), \(2.098\times10^{-14}\)
& YP3-H2 HS\_1\_1 (PASS), \(3.333\times10^{-11}\) & HS (1, 1): the
couplings of h to (d theta\_0)\^{}2 and (d theta\_1)\^{}2 by exact
differentiation of the kinetic part, 0.07096573 and 0.07212812, = f\_F
and f\_I\^{}2/f\_F and = Rust (largest relative difference) \\
WP3-H2 HS\_0\_1 & PASS & \(1.421\times10^{-13}\) &
\(1.000\times10^{-10}\) & P3-H2 HS\_0\_1 (PASS), \(3.419\times10^{-14}\)
& YP3-H2 HS\_0\_1 (PASS), \(3.334\times10^{-11}\) & HS (0, 1): the
couplings of h to (d theta\_0)\^{}2 and (d theta\_1)\^{}2 by exact
differentiation of the kinetic part, 0.067233343 and 0.068220149, = f\_F
and f\_I\^{}2/f\_F and = Rust (largest relative difference) \\
WP3-H2 HS\_0\_0 & PASS & \(7.694\times10^{-14}\) &
\(1.000\times10^{-10}\) & P3-H2 HS\_0\_0 (PASS), \(8.660\times10^{-15}\)
& YP3-H2 HS\_0\_0 (PASS), \(3.333\times10^{-11}\) & HS (0, 0): the
couplings of h to (d theta\_0)\^{}2 and (d theta\_1)\^{}2 by exact
differentiation of the kinetic part, 0.061393961 and 0.06139396, = f\_F
and f\_I\^{}2/f\_F and = Rust (largest relative difference) \\
WP3-H2 KS & PASS & \(1.357\times10^{-13}\) & \(1.000\times10^{-10}\) &
P3-H2 KS (PASS), \(4.763\times10^{-14}\) & YP3-H2 KS (PASS),
\(3.334\times10^{-11}\) & KS: the couplings of h to (d theta\_0)\^{}2
and (d theta\_1)\^{}2 by exact differentiation of the kinetic part,
0.071348166 and 0.060506197, = f\_F and f\_I\^{}2/f\_F and = Rust
(largest relative difference) \\
WP3-H2 CJT & PASS & \(1.352\times10^{-13}\) & \(1.000\times10^{-10}\) &
P3-H2 CJT (PASS), \(7.683\times10^{-14}\) & YP3-H2 CJT (PASS),
\(3.334\times10^{-11}\) & CJT: the couplings of h to (d theta\_0)\^{}2
and (d theta\_1)\^{}2 by exact differentiation of the kinetic part,
0.071348166 and 0.064162461, = f\_F and f\_I\^{}2/f\_F and = Rust
(largest relative difference) \\
WP3-H1 HS\_1\_1 & PASS & -- & -- & P3-H1 HS\_1\_1 (PASS), \(0\) & YP3-H1
HS\_1\_1 (PASS), -- & HS (1, 1): the amplitude mode lies below the
threshold of two quasi-particles: M\_H = 0.8136931 \textless{} 2 E\_gap
= 0.9059348 (the gap of the ground state computed by the notebook)
{[}M\_H / 2 E\_gap = 0.89818{]} \\
WP3-H3 HS\_1\_1 & PASS & \(3.168\times10^{-13}\) &
\(1.000\times10^{-9}\) & P3-H3 HS\_1\_1 (PASS), \(0\) & YP3-H3 HS\_1\_1
(PASS), \(3.168\times10^{-13}\) & HS (1, 1): the numbers of the Higgs
sector (v\_SF = 0.14308913, M\_H/v\_SF = 5.68662, lambda = 16.1688,
Gamma/M\_H = 3.9232) = Rust (largest relative difference) \\
WP3-H1 HS\_0\_1 & PASS & -- & -- & P3-H1 HS\_0\_1 (PASS), \(0\) & YP3-H1
HS\_0\_1 (PASS), -- & HS (0, 1): the amplitude mode lies below the
threshold of two quasi-particles: M\_H = 0.7250834 \textless{} 2 E\_gap
= 0.8271651 (the gap of the ground state computed by the notebook)
{[}M\_H / 2 E\_gap = 0.87659{]} \\
WP3-H3 HS\_0\_1 & PASS & \(1.208\times10^{-13}\) &
\(1.000\times10^{-9}\) & P3-H3 HS\_0\_1 (PASS), \(0\) & YP3-H3 HS\_0\_1
(PASS), \(1.460\times10^{-13}\) & HS (0, 1): the numbers of the Higgs
sector (v\_SF = 0.1354499, M\_H/v\_SF = 5.35315, lambda = 14.3281,
Gamma/M\_H = 3.4708) = Rust (largest relative difference) \\
WP3-H1 HS\_0\_0 & PASS & -- & -- & P3-H1 HS\_0\_0 (PASS), \(0\) & YP3-H1
HS\_0\_0 (PASS), -- & HS (0, 0): the amplitude mode lies below the
threshold of two quasi-particles: M\_H = 0.6143874 \textless{} 2 E\_gap
= 0.7327813 (the gap of the ground state computed by the notebook)
{[}M\_H / 2 E\_gap = 0.83843{]} \\
WP3-H3 HS\_0\_0 & PASS & \(1.850\times10^{-13}\) &
\(1.000\times10^{-9}\) & P3-H3 HS\_0\_0 (PASS), \(0\) & YP3-H3 HS\_0\_0
(PASS), \(1.853\times10^{-13}\) & HS (0, 0): the numbers of the Higgs
sector (v\_SF = 0.12278792, M\_H/v\_SF = 5.00365, lambda = 12.5182,
Gamma/M\_H = 2.9885) = Rust (largest relative difference) \\
WP3-H1 KS & PASS & -- & -- & P3-H1 KS (PASS), \(0\) & YP3-H1 KS (PASS),
-- & KS: the amplitude mode lies below the threshold of two
quasi-particles: M\_H = 0.603476 \textless{} 2 E\_gap = 0.7327813 (the
gap of the ground state computed by the notebook) {[}M\_H / 2 E\_gap =
0.82354{]} \\
WP3-H3 KS & PASS & \(2.253\times10^{-13}\) & \(1.000\times10^{-9}\) &
P3-H3 KS (PASS), \(3.369\times10^{-16}\) & YP3-H3 KS (PASS),
\(2.514\times10^{-13}\) & KS: the numbers of the Higgs sector (v\_SF =
0.13140786, M\_H/v\_SF = 4.59239, lambda = 10.545, Gamma/M\_H = 2.1349)
= Rust (largest relative difference) \\
WP3-H1 CJT & PASS & -- & -- & P3-H1 CJT (PASS), \(0\) & YP3-H1 CJT
(PASS), -- & CJT: the amplitude mode lies below the threshold of two
quasi-particles: M\_H = 0.603476 \textless{} 2 E\_gap = 0.7327813 (the
gap of the ground state computed by the notebook) {[}M\_H / 2 E\_gap =
0.82354{]} \\
WP3-H3 CJT & PASS & \(3.085\times10^{-13}\) & \(1.000\times10^{-9}\) &
P3-H3 CJT (PASS), \(3.573\times10^{-16}\) & YP3-H3 CJT (PASS),
\(3.085\times10^{-13}\) & CJT: the numbers of the Higgs sector (v\_SF =
0.13531998, M\_H/v\_SF = 4.45962, lambda = 9.94412, Gamma/M\_H = 2.1349)
= Rust (largest relative difference) \\
WQN-1 & PASS & -- & -- & QN-1a (PASS), \(0\); QN-1b (PASS), \(0\); QN-1c
(PASS), \(0\); QN-1d (PASS), \(0\); QN-1e (PASS), -- & YQN-1 (PASS), --
& the algebra of SU(2)\_L x SU(2)\_R x U(1)\_R x U(1)\_F\_tot, Hermitian
generators, Y = T\^{}3\_R + R/4 (4.3), Q = T\^{}3 + T\^{}3\_R + R/4, and
{[}Y, T\^{}1\_R{]} != 0 {[}identities of exact matrices{]} \\
WQN-2 & PASS & -- & -- & QN-2a (PASS), --; QN-2b (PASS), \(0\) & YQN-2
(PASS), -- & the six real Lorentz generators: S({[}u, u'{]}) = {[}S(u),
S(u'){]}, independent over the reals, commuting with the ten internal
generators {[}identities of exact matrices; real rank 6{]} \\
WQN-4 & PASS & -- & -- & QN-4a (PASS), \(0\); QN-4b (PASS), \(0\); QN-4c
(PASS), \(0\) & YQN-4 (PASS), -- & Q = 3/4 (up), -1/4 (down) on Psi\_LL,
1/4 on chi\_SH, -3/4 on chi'\_SH; 4 Y = F\_tot on Psi\_LL and chi\_SH,
-3 F\_tot on chi'\_SH; the Casimir operators 3/4 of SU(2)\_L on Psi\_LL
and of SU(2)\_R on the singlets {[}Q on the components 0..15: \{3/4,
-1/4, 3/4, -1/4, 3/4, -1/4, 3/4, -1/4, 1/4, 1/4, 1/4, 1/4, -3/4, -3/4,
-3/4, -3/4\}{]} \\
WQN-5 & PASS & -- & -- & QN-5a (PASS), \(0\); QN-5b (PASS), \(0\); QN-5c
(PASS), -- & YQN-5 (PASS), -- & no anomaly (47 traces exactly 0),
vector-like fields, an even number of Weyl doublets {[}2 Weyl
doublets{]} \\
WQN-6 & PASS & -- & -- & QN-6a (PASS), --; QN-6b (PASS), --; QN-6c
(PASS), --; QN-6d (PASS), \(3.201\times10^{-16}\); QN-6e (PASS), -- &
YQN-6 (PASS), -- & the dimensions 6, 4, 4, 2, 2, 2 (pairs Psi Psi, Psi
chi, Psi chi', chi chi, chi chi', chi' chi'), 8, 4, 4, 2, 2, 2
(particle-antiparticle), 2, 0, 0 (symmetric Psi Psi, chi chi, chi'
chi'), 8, 6, 8, 8 (the pair of singlets as one field); C gamma\^{}ch (x)
e\_s, C (x) e\_s are invariant and span the pairs Psi chi and Psi chi'
{[}dimensions \{6, 4, 4, 2, 2, 2, 8, 4, 4, 2, 2, 2\}, symmetric \{2, 0,
0\}, custodial \{8, 6, 8, 8\}{]} \\
WQN-7 & PASS & -- & -- & QN-7 (PASS), \(0\) & YQN-7 (PASS), \(0\) &
parity: +1 for C gamma\^{}ch (x) e\_s, -1 for C (x) e\_s, with either
singlet {[}identities of exact matrices{]} \\
WQN-8 & PASS & -- & -- & QN-8a (PASS), \(0\); QN-8b (PASS), \(0\); QN-8c
(PASS), \(0\); QN-8d (PASS), \(0\); QN-8e (PASS), \(0\); QN-8f (PASS),
\(0\); QN-8g (PASS), -- & YQN-8 (PASS), -- & on Phi\_SH: tau\^{}a/2, Y =
1/2, Q = diag(1, 0); on Phi'\_SH: tau\^{}a/2, Y = -1/2, Q = diag(0, -1);
F\_tot = 2, R = 0, T\^{}3\_R = +-1/2, Casimir 3/4, closure; i tau\^{}2
Phi'\^{}* has tau\^{}a/2 and Y = 1/2; the bidoublet has (tau\^{}a/2) (x)
1 and 1 (x) (tau\^{}a/2): M\_SH -\textgreater{} L M\_SH R\^{}T {[}G\_Y =
\{\{1/2, 0\}, \{0, 1/2\}\} and \{\{-1/2, 0\}, \{0, -1/2\}\}{]} \\
WQN-9 & PASS & -- & -- & QN-9a (PASS), \(1.056\times10^{-14}\); QN-9b
(PASS), \(1.754\times10^{-14}\); QN-9c (PASS), \(5.251\times10^{-14}\);
QN-9d (PASS), -- & YQN-9a (PASS), \(3.629\times10^{-15}\); YQN-9b
(PASS), \(1.494\times10^{-15}\) & the generators on Phi\_SH, Phi'\_SH
and M\_SH are Hermitian, so the pairing term (4.2) is invariant; det
M\_SH has the charges 0 under T\^{}a, T\^{}a\_R, R and Y and 4 under
F\_tot (the determinant term breaks only U(1)\_F\_tot) {[}charges of det
M\_SH \{0, 0, 0, 0, 0, 0, 0, 0, 4\}{]} \\
WQN-10 & PASS & -- & -- & QN-10a (PASS), \(3.159\times10^{-15}\); QN-10b
(PASS), --; QN-10c (PASS), --; QN-10g (PASS), -- & YQN-10 (PASS), -- &
the 24 channels (parity split, closure, SU(2) multiplets, Y, Q, F\_tot,
R, T\^{}3\_R) and the 8 custodial channels (SU(2)\_L x SU(2)\_R
multiplets) agree with the tables of the Rust program {[}all 24 and all
8 agree; custodial: Psi (chi, chi') (2,2); (chi, chi') (chi, chi')
(1,3); Psi\^{}dag (chi, chi') (2,2); (chi, chi')\^{}dag (chi, chi')
(1,1) + (1,3){]} \\
WQN-10e & PASS & -- & -- & QN-10e (PASS), --; QN-10f (PASS), -- & YQN-10
(PASS), -- & no pair of two particles of one field and no chi chi' pair
is neutral; the pair channels have exactly one neutral component each in
Psi chi and Psi chi' and none elsewhere {[}smallest absolute value of Q
1/2; neutral components of the 12 pair channels \{0, 0, 1, 1, 1, 1, 0,
0, 0, 0, 0, 0\}{]} \\
WQN-11 & PASS & -- & -- & QN-11b (PASS), --; QN-11d (PASS), --; QN-11e
(PASS), -- & YQN-11 (PASS), -- & the stabilizer of (0, phi) and of the
neutral vacuum of H5 in SU(2)\_L x U(1)\_Y is spanned by Q; that of i
tau\^{}2 among the eight generators is spanned by T\^{}a + T\^{}a\_R and
R, three gauged generators and F\_tot are broken; other phases with
equal lengths keep four generators, unequal lengths only T\^{}3 +
T\^{}3\_R and R {[}stabilizer of i tau\^{}2 \{\{0, 0, 0, 0, 0, 0, 1,
0\}, \{0, 0, 1, 0, 0, 1, 0, 0\}, \{0, 1, 0, 0, 1, 0, 0, 0\}, \{1, 0, 0,
1, 0, 0, 0, 0\}\}{]} \\
WQN-11f & PASS & -- & -- & QN-11f (PASS), --; QN-11g (PASS), -- & YQN-11
(PASS), -- & a generic second condensate with a charged component, added
to the neutral vacuum of H5, breaks Q; one in the neutral part of a
channel keeps exactly Q {[}24 of 24 channels as expected{]} \\
WQN-12 & PASS & -- & -- & QN-12b (PASS), \(1.094\times10^{-16}\); QN-12d
(PASS), \(5.551\times10^{-17}\); QN-12f (PASS), \(4.375\times10^{-16}\)
& YQN-12 (PASS), \(4.375\times10^{-16}\) & the mass matrix of (W\^{}1,
W\^{}2, W\^{}3, B) from Phi\_SH = (0, phi) and Phi'\_SH = (phi', 0):
eigenvalues 0, g\^{}2 v\_SH\^{}2/4 (twice), (g\^{}2 + g'\^{}2)
v\_SH\^{}2/4 with v\_SH\^{}2 = 2 (phi\^{}2 + phi'\^{}2) for any phi,
phi', the massless vector along (0, 0, g', g), and rho = 1 (symbolic
couplings and lengths) {[}eigenvalues \{0, (g\^{}2*(ph\^{}2 +
php\^{}2))/2, (g\^{}2*(ph\^{}2 + php\^{}2))/2, ((g\^{}2 +
gp\^{}2)*(ph\^{}2 + php\^{}2))/2\}{]} \\
WSF-2a & PASS & -- & -- & SF-2a (PASS), \(3.508\times10^{-11}\) & YSF-2
(PASS), \(4.270\times10^{-11}\) & the triplet matrix K\^{}3 couples the
up and down components of Psi\_LL by -C gamma\^{}ch in both blocks
(Psi\^{}T K\^{}3 Psi = -2 Psi\_u\^{}T C gamma\^{}ch Psi\_d); K\_d
couples the down components and chi\_SH, K'\_u the up components and
chi'\_SH, by C gamma\^{}ch / 2 in both blocks: the same channel with
half the coefficient, so the pairing term -(G/4) /Psi\^{}T K Psi/\^{}2
has the critical coupling G\^{}crit\_SH = 4 G\^{}crit\_LL {[}identities
of exact matrices{]} \\
WSF-2b & PASS & -- & -- & SF-2a (PASS), \(3.508\times10^{-11}\) & YSF-2
(PASS), \(4.270\times10^{-11}\) & I1(m) = int\_0\^{}1 q\^{}2 /
sqrt(q\^{}2 + m\^{}2) dq = (sqrt(1 + m\^{}2) - m\^{}2 arsinh(1/m)) / 2
(the exact integral equals it at ten masses, to 40 digits);
G\^{}crit\_SH = 4 pi\^{}2 / I1(1/5) = 85.146439022198692144 and G\_SH /
G\^{}crit\_SH = 1.183842814303957 (G\_SH = 504/5), against the Rust
values of both versions {[}deviation 3.508e-11{]} \\
WSF-3 & PASS & -- & -- & SF-3a (PASS), \(7.517\times10^{-12}\); SF-3b
(PASS), -- & YSF-3 (PASS), \(8.011\times10^{-12}\) & the exchange
contraction of the pairing terms: a scalar density n of Psi\_LL gives
each paired singlet the self-energy -(G\_SH/16) n beta and a scalar
density of a paired singlet gives Psi\_LL -(G\_SH/16) n beta (x) 1\_2
(exact); the susceptibility chi\_S = (2/pi\^{}2) int\_0\^{}1
q\^{}4/(q\^{}2 + m\^{}2)\^{}(3/2) dq = 0.08316014730342272 (Psi\_LL;
half for a singlet); the largest eigenvalue of the 3 x 3 scalar response
of CJT for V\^{}(1)\_SH = 0.915274035123 (second version) and
0.789268068370 (first), against the Rust values {[}deviation
1.063e-11{]} \\
WSF-4 & PASS & -- & -- & SF-4a (PASS), \(1.350\times10^{-11}\) & none &
the scalar triplet pair coupling of KS and CJT, -(1/4) G\^{}S - G\^{}V =
6/5, and its eigenvalue in the normal vacuum, (6/5) I1(1/5) / pi\^{}2 =
0.05637346734780750, against the Rust value {[}deviation 8.097e-12{]} \\
WSF-5 & PASS & -- & -- & SF-6a (PASS), --; SF-6b (PASS), -- & YSF-6
(PASS), \(4.955\times10^{-10}\); YSF-12 (PASS), \(1.562\times10^{-12}\)
& HS (0, 1) and HS (0, 0), where only the pairing terms act: the gap
equation G\_SH I1(M\_eff) / (4 pi\^{}2) = 1 gives M\_eff =
0.4529674056079986 and /D/ = 2 (M\_eff\^{}2 - m\^{}2)\^{}(1/2) =
0.8128455463204339 (30 digits); the singular values of the bidoublet of
the Rust program are /D/ and /D/ in the second version and /D/ and 0 in
the first {[}deviation 4.841e-12{]} \\
WH-1 & PASS & -- & -- & H-4f (PASS), \(2.447\times10^{-13}\); H-4g
(PASS), \(8.096\times10^{-13}\); H-5a (PASS), \(1.703\times10^{-13}\);
H-5b (PASS), \(2.018\times10^{-13}\); H-6b (PASS),
\(1.569\times10^{-13}\) & YH-2 (PASS), \(2.882\times10^{-9}\) & the
effective Lagrangian Z tr((dM)\^{}dag dM) of the bidoublet at M = phi i
tau\^{}2: chi(T\^{}a, T\^{}b) = chi(T\^{}a\_R, T\^{}b\_R) = -chi(T\^{}a,
T\^{}b\_R) = Z phi\^{}2 delta\_ab; with Y = T\^{}3\_R + R/4 the charge Q
= T\^{}3 + Y does not respond; the mass matrix of (W\^{}1, W\^{}2,
W\^{}3, B) has the eigenvalues 0, g\^{}2 Z phi\^{}2 (twice), (g\^{}2 +
g'\^{}2) Z phi\^{}2, so M\_W = g v/2 and M\_Z = (g\^{}2 +
g'\^{}2)\^{}(1/2) v/2 with v\^{}2 = 4 Z phi\^{}2, and rho = 1 (symbolic)
{[}eigenvalues \{0, g\^{}2*phi\^{}2*zz, g\^{}2*phi\^{}2*zz, (g\^{}2 +
gp\^{}2)*phi\^{}2*zz\}{]} \\
WH-2 & PASS & -- & -- & none (the column M\_EFT) & YH-5 (PASS),
\(1.768\times10^{-16}\) & the effective Lagrangian of the pair fields
gives the mass of the amplitude mode M\^{}2 = V'\,' phi\_0\^{}2 / (4
chi\^{}\{33\}) in both versions: second version chi\^{}\{33\} =
phi0\^{}2*zz, M\^{}2 = V'`/(4 Z); first version chi\^{}\{33\} =
(phi0\^{}2*zz)/2, M\^{}2 = V''/(2 Z) (symbolic) {[}identities{]} \\
WH-3 & PASS & -- & -- & H-2b (PASS), --; H-2c (PASS),
\(2.302\times10^{-4}\); H-8a (PASS), \(8.549\times10^{-14}\) & YH-1
(PASS), \(6.455\times10^{-13}\); YH-3 (PASS), \(2.016\times10^{-10}\) &
HS (0, 0): the gap equation (G\_SH / 4 pi\^{}2) I1(M\_eff) = 1 gives
M\_eff = 0.4529674056079986, d\_0 = 0.8128455463204339, phi\_0 = d\_0 I1
/ (4 pi\^{}2) = d\_0 / G\_SH = 0.008063943911909066; dI1/dM = -M J with
J = arsinh(1/M) - (1 + M\^{}2)\^{}(-1/2) (exact; J equals its integral
to 40 digits); V'\,'(phi\_0) = 2 pi\^{}2 d\_0\^{}2 J / (I1 (I1 -
d\_0\^{}2 J/4)) = 71.68031758679552 for one pair (first version) and
twice that, 143.3606351735910, for the two pairs of the second version;
against the Rust values (relative) {[}deviation 1.260e-8{]} \\
WH-4 & PASS & -- & -- & none (the columns in GeV) & YH-5 (PASS),
\(1.768\times10^{-16}\) & m\_H / v = 125.13 / v = 0.5082047224 (the
inputs of stage 2); for every row of the masses, Lambda\_LL = v / v\_SH,
M\_H = omega\_H Lambda\_LL and M\_H / v = omega\_H / v\_SH (relative);
the amplitude modes of the linear response give M\_H / v from 7.96732 to
10.1624, 15.68 to 20.00 times the value of the Standard Model
{[}deviation 8.766e-17{]} \\
WH-5 & PASS & -- & -- & H-7d (PASS), \(1.990\times10^{-13}\); H-8d
(PASS), \(1.030\times10^{-12}\) & YH-7 (PASS), \(9.853\times10^{-16}\) &
HS (0, 0): the identity (k\^{}2 + m\^{}2) / (E (E\^{}2 - w\^{}2/4)) -
1/E = (w\^{}2 - d\^{}2) / (4 E (E\^{}2 - w\^{}2/4)) with E\^{}2 = k\^{}2
+ m\^{}2 + d\^{}2/4 holds for every momentum (symbolic) and E\^{}2 -
w\^{}2/4 \textgreater{} 0 for w \textless{} 2 M\_eff (symbolic), so that
lambda\_amp(w) = 1 only at w = d\_0; the continuum integrals give
/lambda\_amp(d\_0) - 1/ = 0 and /lambda\_phase(0) - 1/ = 0 (the gap
equation), the root of lambda\_amp = 1 is 0.81284554632043386736 against
d\_0 = 0.81284554632043386736, below 2 M\_eff = 0.9059348112159973
(binding 0.0930893); the amplitude modes of the Rust program (HS (0, 1)
and HS (0, 0), both versions) equal d\_0 (relative) {[}identity True,
positivity True; Rust deviation 8.442e-13{]} \\
WD-1 & PASS & -- & -- & D-3a (PASS), \(1.097\times10^{-6}\); D-3b
(PASS), \(6.115\times10^{-10}\); D-3c (PASS), \(1.991\times10^{-8}\) &
YD-2a (PASS), \(3.611\times10^{-7}\); YD-2b (PASS),
\(5.922\times10^{-9}\) & the time-dependent mean field i dR/dt = {[}H,
R{]} conserves tr(H R) when H is the derivative of the energy (d tr(H
R)/dt = -i tr(H {[}H, R{]}) = 0), keeps R Hermitian, and keeps a
projector a projector (symbolic, 3 x 3 Hermitian matrices with symbolic
entries; the projector of rank 1 along the first axis) {[}energy True,
hermiticity True, projector True{]} \\
WD-2 & PASS & -- & -- & D-6a (PASS), \(1.045\times10^{-9}\); D-6b
(PASS), \(1.481\times10^{-9}\) & YD-5 (PASS), \(1.045\times10^{-9}\) &
the Josephson relation: exp(i mu F t) R\_0 exp(-i mu F t) solves i dR/dt
= {[}H - mu F, R{]} for a stationary R\_0 and multiplies the pair block
by exp(2 i mu t) (symbolic); the phases of d(t) of the Josephson runs of
the Rust program grow at the rate 2 mu = 0.02 (straight-line fit in
exact arithmetic on the numbers of the file; relative) {[}rates
josephson\_second\_HS0\_0 0.0199999999791, josephson\_second\_CJT
0.0200000000125; deviation 1.045e-9{]} \\
WD-3 & PASS & -- & -- & D-4a (PASS), \(1.239\times10^{-6}\) & YD-3
(PASS), \(1.241\times10^{-6}\) & the peaks of the amplitude runs of HS
(0, 1) and HS (0, 0) of the Rust program (both versions) lie at d\_0 =
0.81284554632043386736 of the gap equation (relative) {[}deviation
1.239e-6{]} \\
WD-4 & PASS & -- & -- & D-4b (PASS), \(5.048\times10^{-6}\) & none & the
Hann window: its normalized transform has the modulus /sin(x pi) / (x pi
(1 - x\^{}2))/ at the offset x 2 pi / T (the squared modulus of the
exact integral equals that of the closed form at nine rational offsets,
exactly) and the value 1/2 at x = 1 (limit), so that the half width at
half maximum of a pure cosine is 2 pi / T = 0.020943951 for the runs of
length 300; the widths of the sampled cosines of the Rust program
against it (relative) {[}form True, half True, deviation 5.037e-7{]} \\
WD-5 & PASS & -- & -- & none & YD-6 (PASS), \(6.022\times10^{-9}\) & HS
(0, 0): the reduced model of 48 Anderson pseudospins (one at +eps\_k and
one at -eps\_k for each of the 24 radial nodes, eps\_k = (k\^{}2 +
m\^{}2)\^{}(1/2)), integrated here by NDSolve with the method IDA, gives
d\_0 = 2 Delta\_0 = 0.812845546320434 from its gap equation and, after
the kick 1e-3, the d(t) of the Rust program for 0 \textless= t
\textless= 50 (largest /difference/ / d\_0, at the 51 stored times)
{[}d\_0 against the gap equation of WD-3 6.032e-22; deviation
1.749e-9{]} \\
\end{longtable}

\textbf{Table 1.} The checks of the notebook. Data:
\texttt{data/hck\_notebook\_log.txt},
\texttt{data/hck\_part1\_checks.txt},
\texttt{data/hck\_part1\_checks\_python.txt},
\texttt{data/hck\_part2\_checks.txt},
\texttt{data/hck\_part2\_checks\_python.txt},
\texttt{data/hck\_part3\_checks.txt},
\texttt{data/hck\_part3\_checks\_python.txt}, and for Part R
\texttt{data/hck\_sh\_quantum\_checks.txt},
\texttt{data/hck\_sh\_superfluid\_checks.txt},
\texttt{data/hck\_sh\_higgs\_checks.txt},
\texttt{data/hck\_sh\_dynamics\_checks.txt} and their Python reports
(\texttt{data/hck\_sh\_quantum\_checks\_python.txt} and so on).

\hypertarget{the-methods-of-the-notebook-and-of-its-counterparts}{%
\section{7. The methods of the notebook and of its
counterparts}\label{the-methods-of-the-notebook-and-of-its-counterparts}}

The three implementations compute the same quantities by different
means, so that an error of one method does not repeat itself in the
others:

\begin{itemize}
\item
  \textbf{bi-quaternions and bi-octonions} (WA-2 to WA-12): the notebook
  uses Gaussian-integer coefficients and exact arithmetic, so every
  identity holds with the deviation exactly 0; the Rust and the Python
  programs use small integer coefficients in double precision (also
  exact) and random floating-point values where the notebook uses exact
  ones;
\item
  \textbf{Lorentz and spinor matrices} (WA-10, WF-3): floating point in
  all three, with different random group elements (the seeds of the
  three generators differ);
\item
  \textbf{theorem D1} (WF-5): exact ranks of exact matrices in the
  notebook; ranks by the singular values in Python, and by the Rust
  program's own elimination;
\item
  \textbf{covariant derivatives} (WC-1 to WC-8): symbolic derivatives
  (the function \texttt{D}) of the random smooth fields in the notebook;
  automatic differentiation in Rust; central differences of eighth order
  in Python;
\item
  \textbf{the field equations as commutators} (WE-2 to WE-4): exact
  operator identities on the 256 states of one point in the notebook
  (exact couplings \(-8\), \(4/5\), \(45/2\), deviation exactly 0);
  double precision in the two programs;
\item
  \textbf{the Yang-Mills equations} (WE-14): the function
  \texttt{EulerEquations} of the package \texttt{VariationalMethods} in
  the notebook; the program sympy in Python; in Rust, exact derivatives
  of the Lagrangian density, a polynomial of degree four in the fields
  and their first derivatives (five-point differences in its arguments,
  exact for such a polynomial, and automatic differentiation along the
  coordinates);
\item
  \textbf{the conventions of the Standard Model} (WSM-1): exact, with
  symbolic couplings, in the notebook (the characteristic polynomial and
  the null vector of the mass matrix); the four eigenvectors with 20
  values of the couplings in Rust; the eigenvalues by LAPACK with 20
  random values in Python;
\item
  \textbf{the analytic solutions} (WAS1-3 to WAS4-3): dense eigenvalues
  of the Hermitian matrices by the function \texttt{Eigenvalues} (up to
  4800 rows), the gap equation by \texttt{FindRoot} with 30 digits, and
  the profiles of the kink by \texttt{FullSimplify} in the notebook; the
  Rust program's own eigensolvers (inertia counting and bisection of
  block-tridiagonal matrices) and the band solver of LAPACK in Python;
\item
  \textbf{the pair channels} (WP2-0a, WP2-0b): exact operator algebra on
  the 256 states of one point with exact matrices in the notebook
  (deviation exactly 0) and in floating point in Python; the Rust
  program checks the same table against the pair contraction of the mean
  field (P2-0e) and the sum rule of the three W (P2-0d);
\item
  \textbf{the engine of the three methods} (WP2-1 to WP2-8d): the
  notebook ports the formulas of the energy and of the self-energies,
  the grid (Gauss-Legendre nodes from \texttt{NSolve} of the Legendre
  polynomial with 40 digits) and the iteration, and uses
  \texttt{Eigensystem} of the exactly Hermitian (symmetrized)
  \(16\times16\) matrices, \texttt{Eigenvalues} of the \(120\times120\)
  response matrices, \texttt{LinearSolve} for \(V^{(2)}_{\rm SF}\), and
  \texttt{NIntegrate} with 30 digits for the angular averages of the
  Yukawa kernel; the Rust program uses its own eigensolvers (Householder
  and QL for Hermitian matrices, the shifted QR algorithm for general
  ones) and the Python program LAPACK. The three programs share the grid
  (24 radial nodes, six directions), the starts and the tolerances, so
  that the same states are found; the random states of WP2-1 differ;
\item
  \textbf{the accuracy of the Landau form} (WP3-8, WP3-10, WP3-11): the
  least squares by the function \texttt{LeastSquares} on scaled columns
  in the notebook, by Householder reflections in Rust, and by the
  singular-value decomposition of numpy in Python; the quasi-particle
  energies by \texttt{Eigenvalues} of the \(16\times16\) matrices (the
  Rust program's own eigensolver, LAPACK in Python); the Taylor
  coefficient at the origin from the first six points of the computed
  curve in the notebook, and from three states with small applied fields
  and the Richardson extrapolation in the two programs; the integral
  with the term \(\tfrac1{16}s^2\ln s\) exactly, by \texttt{Integrate},
  \texttt{FullSimplify} and \texttt{Series}, in the notebook alone (it
  is shown, not checked);
\item
  \textbf{the free energy at finite temperature} (WP3-T1 to WP3-T9): the
  thermal states \(f(\mathsf H_{\rm SF})\) by \texttt{Eigensystem} of
  the symmetrized \(16\times16\) matrices on the grid of Part II in the
  notebook, by the Rust program's own eigensolver and by LAPACK in
  Python; the critical temperatures by bisection, in the notebook with
  the Dirac mass of the normal state from its scalar self-consistency
  alone (and, for the three methods with a scalar coupling, compared
  with the normal state of the thermal engine), in Rust and Python with
  the normal state of the engine (and, in Python's YP3-T9, also from the
  scalar self-consistency); the Taylor coefficients by
  \texttt{LinearSolve} of the cubic interpolation in the scaled variable
  \(\Phi^2/\Phi^2_{\max}\) (Gaussian elimination in Rust,
  \texttt{numpy.linalg.solve} in Python); the residuals of the Landau
  forms of the Rust free-energy curves by \texttt{LeastSquares}
  (Householder reflections in Rust, the singular-value decomposition in
  Python);
\item
  \textbf{the width of the critical region and the Higgs sector} (WP3-G1
  to WP3-G12, WP3-H1 to WP3-H3): the static pair response by
  \texttt{Eigensystem} of the symmetrized \(8\times8\) particle and hole
  blocks at every node, the matrices
  \(\mathsf U_A^\dagger(E_{ab}-E_{ba})\mathsf U_D\) from the components
  of the eigenvectors by \texttt{Outer}, and the system of the other 25
  pair channels by \texttt{LinearSolve} in the notebook; the Rust
  program's own eigensolver and LU decomposition, and LAPACK (batched
  \texttt{eigh}, \texttt{solve}) in Python. The notebook and the Rust
  program work on the grid of Part II (24 radial nodes, six directions),
  the Python program on 64 directions; the product rule of the two
  regularizations (24 radial nodes, 24 nodes in \(\cos\theta\), 6 angles
  \(\varphi\)) is the same in the three. The slope in closed form uses
  the derivatives of \(I^{(1)}_{\rm SF}\) on the radial nodes of the
  grid in the notebook and in Rust, on 400 Gauss-Legendre nodes in
  Python. The couplings of the amplitude field come from exact symbolic
  differentiation of the kinetic part (\texttt{D},
  \texttt{ComplexExpand}, \texttt{Simplify}) in the notebook, from
  finite differences of the field in Python and of the kinetic part in
  Rust. The twisted states of the nonlinear engine of HS \((0,0)\)
  (P3-G7) are not repeated by the notebook.
\item
  \textbf{the revised hypotheses} (Part R, WQN-1 to WD-5): exact
  arithmetic and symbolic algebra in the notebook, as in the Wolfram
  scripts of the revision whose code it repeats (the generators as exact
  rational matrices, the invariant bilinears and the stabilizers as
  exact null spaces, the identities of the effective Lagrangian and of
  the time-dependent mean field with symbolic entries, the integrals by
  \texttt{Integrate} or by \texttt{NIntegrate} with 30 to 40 digits, the
  positivity of the factor of remark R11 by \texttt{Resolve} over the
  real numbers, and the reduced model of stage 6 by \texttt{NDSolve}
  with the method IDA); the Rust programs of the revision compute the
  same quantities in double precision with their own eigensolvers and,
  for the dynamics, with CVODE of SUNDIALS on all sixteen components;
  the Python programs with numpy and LAPACK and, for the dynamics, with
  the same CVODE through its C interface.
\end{itemize}

\hypertarget{the-symbols-of-the-code-of-part-i}{%
\section{8. The symbols of the code of Part
I}\label{the-symbols-of-the-code-of-part-i}}

Every name of the code that stands for an object of the physics is
introduced in the notebook with a comment that gives its printed form
and its meaning. The table lists them in the order of the notebook. (The
names \texttt{dA2}, \texttt{dF3a}, \texttt{dC1}, \texttt{dE2} and their
like hold the measured deviation of the check whose identifier they
contain; the names that end in \texttt{T}, such as \texttt{uT} or
\texttt{psiT}, are the random test fields of Section I.4.)

\begin{longtable}[]{@{}
  >{\raggedright\arraybackslash}p{(\columnwidth - 4\tabcolsep) * \real{0.3333}}
  >{\raggedright\arraybackslash}p{(\columnwidth - 4\tabcolsep) * \real{0.3333}}
  >{\raggedright\arraybackslash}p{(\columnwidth - 4\tabcolsep) * \real{0.3333}}@{}}
\toprule\noalign{}
\begin{minipage}[b]{\linewidth}\raggedright
name in the code
\end{minipage} & \begin{minipage}[b]{\linewidth}\raggedright
printed form
\end{minipage} & \begin{minipage}[b]{\linewidth}\raggedright
meaning
\end{minipage} \\
\midrule\noalign{}
\endhead
\bottomrule\noalign{}
\endlastfoot
\texttt{check}, \texttt{checkTrue}, \texttt{sci}, \texttt{checkHook} &
-- & the functions that record and print a check, the number format of
the reports, and the hook used by the dry run \\
\texttt{mLL} & \(m_{\rm LL}\) & the mass of the spinor field, \(1/5\) in
units of the cutoff \\
\texttt{sig{[}{[}p+1{]}{]}}, \texttt{sigbar{[}{[}p+1{]}{]}} &
\(\sigma^{\rm LL}_p\), \(\bar\sigma_p\) & the unit matrix and the Pauli
matrices; \(\sigma^{\rm LL}_0\) and \(-\sigma^{\rm LL}_j\) \\
\texttt{eLL}, \texttt{etaLL}, \texttt{levi} & \(\mathsf E_{\rm LL}\),
\(\eta^{\rm LL}_{pp'}\), \(\epsilon^{\rm LL}\) & \(i\sigma^{\rm LL}_2\),
the Minkowski metric, the Levi-Civita symbol \\
\texttt{gam{[}p{]}}, \texttt{gch}, \texttt{cLL}, \texttt{betaLL},
\texttt{alphaLL{[}j{]}} & \(\gamma^p_{\rm LL}\),
\(\gamma^{\rm ch}_{\rm LL}\), \(C_{\rm LL}\), \(\beta_{\rm LL}\),
\(\alpha^j_{\rm LL}\) & the Dirac matrices of the chiral basis,
chirality, charge conjugation, and the matrices of the Dirac
Hamiltonian \\
\texttt{d8}, \texttt{i8} & \(A\otimes\mathsf 1_2\),
\(\mathsf 1_4\otimes t\) & a Dirac matrix and an isospin matrix acting
on the eight components \\
\texttt{tLL{[}iota{]}}, \texttt{kLL{[}iota{]}}, \texttt{hLL{[}q,\ m{]}}
& \(T^\iota_{\rm LL}\), \(\mathsf K^\iota_{\rm LL}\),
\(\mathsf h_{\rm LL}(\vec q_{\rm LL})\) & the isospin generators, the
matrices of the scalar spin-zero pairs, the free Dirac Hamiltonian \\
\texttt{qMul}, \texttt{qBar}, \texttt{qDag}, \texttt{qNorm} &
\(\mathbf b\mathbf b'\), \(\bar{\mathbf b}\), \(\mathbf b^\dagger\),
\(N_{\rm LL}(\mathbf b)\) & the product, the conjugations and the
complex norm of bi-quaternions \\
\texttt{mLLmat}, \texttt{biqOf} & \(M_{\rm LL}(\mathbf b)\) & the matrix
of a bi-quaternion and the inverse map \\
\texttt{oMul}, \texttt{oUnit{[}p{]}}, \texttt{oBar}, \texttt{oNorm},
\texttt{oMax} & \(\mathbf x\mathbf y\), \(\mathbf e_p\),
\(\bar{\mathbf x}\), \(N_{\rm LL}(\mathbf x)\) & the product of Cayley
and Dickson, the units, the conjugation and the norm of bi-octonions,
and the size of a difference \\
\texttt{expTraceless}, \texttt{lorentzGenerator},
\texttt{lorentzMatrix}, \texttt{sSpin}, \texttt{sGen},
\texttt{randUnitBiq} & \(\exp\), \(M_{\rm LL}(\mathbf u)\),
\(\mathsf L_{\rm LL}(\mathbf g)\), \(\mathsf S_{\rm LL}(\mathbf g)\),
\(\mathsf s_{\rm LL}(\mathbf u)\) & the exponential of a traceless
\(2\times2\) matrix, a Lorentz generator, the Lorentz matrix, the spinor
matrix, the action of a generator on the Dirac index, a random unit
bi-quaternion \\
\texttt{fockOps}, \texttt{cOps}, \texttt{cdOps} & \(\Psi_{{\rm LL},c}\),
\(\Psi^\dagger_{{\rm LL},c}\) & the annihilation and creation operators
of the eight components of one point (256 states) \\
\texttt{gSM}, \texttt{gpSM}, \texttt{vSM}, \texttt{phiSM}, \texttt{xSM},
\texttt{mSM}, \texttt{photonSM} & \(g\), \(g'\), \(v\),
\(\langle\Phi^{\rm SM}\rangle\), -- , -- , \(A_\mu\) & the couplings and
the vacuum value of the Standard Model, the vacuum, the vectors \(X_a\)
of \(D_\mu\langle\Phi^{\rm SM}\rangle=\sum_aV^a_\mu X_a\), the mass
matrix of \((W^1,W^2,W^3,B)\), the photon \\
\texttt{lorGens}, \texttt{pgGens}, \texttt{amatD1}, \texttt{basisD1},
\texttt{invDim} & -- & theorem D1: the Lorentz and the Pauli-Gürsey
generators on \((\psi,\psi^\dagger)\), the map to the doublet of Weyl
spinors, the antisymmetric bilinears, the dimension of the invariant
ones \\
\texttt{xs}, \texttt{gvT}, \texttt{gfT} & \(x^p_{\rm LL}\),
\(g^{\rm v}_{\rm LL}\), \(g^{\rm F}_{\rm LL}\) & the coordinates as
symbols and the couplings of the test, \(0.83\) and \(0.47\) \\
\texttt{uT}, \texttt{thT}, \texttt{wT{[}p{]}}, \texttt{bT{[}p{]}},
\texttt{wTp{[}p{]}}, \texttt{bTp{[}p{]}} & \(\mathsf U_{\rm LL}\),
\(\theta^{\rm F}_{\rm LL}\), \(\mathsf W_{{\rm SF},p}\),
\(B_{{\rm SF},p}\), \(\mathsf W'_{{\rm SF},p}\), \(B'_{{\rm SF},p}\) &
the test fields of the covariance checks and their gauge transforms \\
\texttt{psiT}, \texttt{psiTp}, \texttt{covPsi}, \texttt{phiT},
\texttt{phiTp}, \texttt{rotT}, \texttt{covPhi}, \texttt{fieldStrength} &
\(\Psi_{\rm LL}\), \(\nabla^{\rm GEE}_p\Psi_{\rm LL}\), \(\vec\Phi\),
\(\mathsf R^{\rm v}_{\rm LL}\), \(\nabla_p\vec\Phi\),
\(\mathsf W_{{\rm SF},pp'}\) & the test spinor and its transform, the
covariant derivative, the isotriplet and its rotation, its covariant
derivative, the field strength \\
\texttt{mF}, \texttt{sF}, \texttt{sFinv}, \texttt{psi4} &
\(M_{\rm LL}(\mathbf g(x))\), \(\mathsf S_{\rm LL}(\mathbf g(x))\) & the
field of unit bi-quaternions of the local \(SL(2,\mathbb C)\) check, its
spinor matrix and inverse, and the test Dirac spinor \\
\texttt{gScal}, \texttt{gVec}, \texttt{gPair}, \texttt{hS}, \texttt{hV},
\texttt{hP}, \texttt{piOps} & \(G^{\rm S}_{\rm SF}\),
\(G^{\rm V}_{\rm SF}\), \(G_{\rm LL}\), -- , -- , -- ,
\(\Pi^\iota_{\rm LL}\) & the exact couplings \(-8\), \(4/5\), \(45/2\)
of the operator checks, the Hamiltonians of the three contact terms of
one point, the pair operators \\
\texttt{wF}, \texttt{jF}, \texttt{fYM}, \texttt{lYM}, \texttt{eqsYM},
\texttt{gYM} & \(W^\iota_{{\rm SF},p}\), \(j^{\iota,p}_{\rm SF}\),
\(W^\iota_{{\rm SF},pp'}\), -- , -- , \(g^{\rm v}_{\rm LL}\) & the
symbolic gauge fields and currents, the field strengths, the Lagrangian
density of the Yang-Mills check, its Euler-Lagrange equations, the
coupling \\
\texttt{hPW}, \texttt{i1}, \texttt{gapAS3} & -- , \(I_1\), \(\Delta\) &
the plane-wave matrix of AS1, the integral of the gap equation of AS3,
the gap at \(G=22.5\) \\
\texttt{kPs{[}iota{]}}, \texttt{pairState}, \texttt{channelOps},
\texttt{channelMats}, \texttt{channelStates}, \texttt{pairEnergy},
\texttt{rustPairEnergy} & \(\mathsf K'^\iota_{\rm LL}\),
\(\lvert\mathsf K\rangle\), \(E^{\rm pair}_{\rm SF}\) & Part II: the
pseudoscalar pair matrices, the state of two charges in a channel, the
seven contact operators, the six channels and their states, the pair
energies and the table of the Rust program \\
\texttt{rustPart2}, \texttt{rustPairChannels} & -- & Part II: the values
of the Rust program \texttt{sf\_part2} and its table of pair-channel
energies \\
\texttt{lamLL}, \texttt{gvLL}, \texttt{gfLL}, \texttt{id8}, \texttt{b8},
\texttt{k8{[}iota{]}}, \texttt{zero8} & \(\Lambda_{\rm LL}\),
\(g^{\rm v}_{\rm LL}\), \(g^{\rm F}_{\rm LL}\), \(\mathsf 1_8\),
\(\beta_{\rm LL}\otimes\mathsf 1_2\), \(\mathsf K^\iota_{\rm LL}\) &
Part II: the cutoff, the gauge couplings 0.7 and 0.45, and the numerical
matrices of the engine \\
\texttt{gaussLegendre}, \texttt{gridSF}, \texttt{ksProj},
\texttt{terms}, \texttt{interaction} & \(\mathcal P^{\rm KS}_{\rm SF}\)
& Part II: the Gauss-Legendre rule, the grid of momenta and weights, the
projection of Kohn and Sham, the contact terms of a method, and the
energy and self-energies of the contact terms \\
\texttt{makeModel}, \texttt{bdg}, \texttt{stateOfNode}, \texttt{states},
\texttt{localDens}, \texttt{solveSF} &
\(\mathsf H_{\rm SF}(\vec q_{\rm LL})\), \(\mathsf N_{\rm SF}\),
\(\mathsf P_{\rm SF}\) & Part II: the free matrices and the normal
vacuum on the grid, the matrices of Bogoliubov and de Gennes, the ground
state of one and of all nodes, the local matrices, and the
self-consistent state \\
\texttt{qpSF}, \texttt{gapSF}, \texttt{dGap}, \texttt{diracMass},
\texttt{effCouplings}, \texttt{methodsSF}, \texttt{shortName} &
\(E^{X,{\rm gap}}_{\rm SF}\), \(\lvert D^X_{\rm SF}\rvert\),
\(M^{X,{\rm D}}_{\rm SF}\), \(g^{X,{\rm S}}_{\rm SF}\),
\(g^{X,{\rm P}}_{\rm SF}\) & Part II: the smallest quasi-particle energy
at a momentum, the gap, the pair amplitude, the Dirac mass, the
effective couplings, the five methods and their names \\
\texttt{dP21}, \texttt{moAS3}, \texttt{solAS3}, \texttt{d3AS3},
\texttt{gCrit}, \texttt{gS2}, \texttt{gP2}, \texttt{potentialsSF} &
\(G^{\rm crit}_{\rm LL}\), \(G^{\rm S}_{\rm SF}\), \(G_{\rm LL}\) & Part
II: the deviation of WP2-1, the pairing-only model of AS3 and its
solution, the critical pair coupling, the coefficients of
\(V^{(2)}_{\rm SF}\), and the three potentials \\
\texttt{devSF}, \texttt{groundSF}, \texttt{groundRows} & -- & Part II:
the deviations from the Rust values, the ground states of the 15 cases,
and the printed table \\
\texttt{pairsSF}, \texttt{toVec}, \texttt{fromVec}, \texttt{respKernel},
\texttt{respFields}, \texttt{respExtract}, \texttt{maskA},
\texttt{maskB}, \texttt{response}, \texttt{mMat}, \texttt{countModes},
\texttt{modeRows} & \(\mathsf F\), \(\mathsf A_n\), \(\mathsf B_n\),
\(\Omega_n\), \(\mathsf M(\omega)\) & Part II: the 120 amplitudes of a
change of the self-energies and their maps, the kernel of the contact
terms, the 120 fields, the response matrices of the nodes,
\(\mathsf M(\omega)=\chi(\omega)\mathsf F\), the number of its real
eigenvalues above 1, and the table of the modes \\
\texttt{chargesSF}, \texttt{susceptibilitySF}, \texttt{chiRows} &
\(\chi^{ab}_{\rm SF}\) & Part II: the four charges, the static
susceptibilities and their table with the masses \(M_W\), \(M_Z\) \\
\texttt{kernSF}, \texttt{kernIntegral}, \texttt{verticesSF},
\texttt{makeBosons}, \texttt{exchangeSF}, \texttt{thoulessSF},
\texttt{cjtExchange} & \(D_0\), \(D_1\) & Part II: the radial kernels of
the angular averages, their integrals with the subtracted logarithm, the
vertices and the four bosons, the exchange self-energies and energy, the
Thouless eigenvalues, and the CJT state with the exchange \\
\texttt{lfVec}, \texttt{lfNorm2}, \texttt{lfIdentity}, \texttt{lfPolar}
& \(\vec\Phi\), \(\lvert\vec\Phi\rvert^2\) & Part III: a complex
3-vector with symbolic real and imaginary parts (\texttt{lfRe{[}k{]}},
\texttt{lfIm{[}k{]}}), its squared norm, the identity
\(\lvert\vec\Phi\rvert^4-\lvert\vec\Phi\cdot\vec\Phi\rvert^2=\lvert\vec\Phi\times\vec\Phi^*\rvert^2\)
and its value on a polar field (\texttt{lfTh}, \texttt{lfR},
\texttt{lfT1}, \texttt{lfT2}: its phase, modulus and direction) \\
\texttt{rustPart3}, \texttt{rustEffPot{[}tag{]}} & -- & Part III: the
values of the Rust program \texttt{sf\_part3} and its five computed
polar curves (applied field, field, energy per volume, Dirac mass) \\
\texttt{lfTags}, \texttt{lfName}, \texttt{lfPhi0}, \texttt{lfE0},
\texttt{lfCurve}, \texttt{lfFitted} & \(X\), \(\Phi^X_{0,{\rm SF}}\),
\(\mathcal V^X_{\rm SF}(\Phi^X_{0,{\rm SF}})\) & Part III: the five
methods, their names, the vacuum field, the energy of the minimum (minus
the depth), the polar curve \(\{\Phi,\mathcal V^X_{\rm SF}\}\), and its
points with \(\Phi\le\) a given multiple of \(\Phi^X_{0,{\rm SF}}\) \\
\texttt{evenFit}, \texttt{landauMin}, \texttt{lfResults},
\texttt{lfKeys}, \texttt{lfRows}, \texttt{pct} & \(\sum_kc_k\Phi^{2k}\),
\(\Phi_{\rm L}\) & Part III: the least-squares fit with the powers
\(\Phi^2,\dots,\Phi^{2K}\), the minimum of the Landau polynomial with
its value and second derivative, all numbers of the accuracy for each
method, the names compared with the Rust values, the rows of the printed
table, a number in per cent \\
\texttt{lfColors}, \texttt{lfShown}, \texttt{lfLegend} & -- & Part III:
the colors, the four distinct curves (KS and CJT coincide) and the
legend of the figure \\
\texttt{lfBdG}, \texttt{lfQpDev} &
\(\mathsf H_{\rm SF}(\vec q_{\rm LL})\) & Part III: the matrix of
Bogoliubov and de Gennes with a Dirac mass and an applied pair field,
and the deviation of its eigenvalues from \(\pm E\) \\
\texttt{lfq}, \texttt{lfs}, \texttt{lfLam}, \texttt{lfInt},
\texttt{lfClosed}, \texttt{lfSeries} & \(q\), \(s\),
\(\Lambda_{\rm LL}\) & Part III: the variables of the integral
\(\int_0^{\Lambda}q^2(q^2+s)^{1/2}dq\), its value, its closed form and
its expansion about \(s=0\) \\
\texttt{lfTaylor} & \(a^{X,{\rm Tay}}_{\rm SF}\) & Part III: the Taylor
coefficient of \(\lvert\vec\Phi\rvert^2\) at the origin, estimated from
the first six points of the curve \\
\texttt{fermiSF}, \texttt{binEntropy}, \texttt{thermalNode},
\texttt{solveSFT}, \texttt{fixedPairT}, \texttt{phiOf} &
\(f(\varepsilon)\), \(\mathcal S_{\rm SF}\), \(\Phi^X_{0,{\rm SF}}(T)\)
& Section III.T: the Fermi function at the temperature \(T\), the
entropy of one mode, the thermal state \(f(\mathsf H_{\rm SF})\) of one
node with the entropy of its 16 modes, the self-consistent thermal
state, the thermal state of an applied pair field, and the order
parameter of a state \\
\texttt{rustThermal}, \texttt{rustThermalLandau},
\texttt{rustThermalCurves}, \texttt{thCol}, \texttt{thRow},
\texttt{thGet}, \texttt{thTaus} & -- & Section III.T: the thermal
results of the Rust program (the table of the order parameter, the
residuals of the Landau forms, the free-energy curves), the names of the
columns of the table, one row and one entry, and the reduced
temperatures of the table \\
\texttt{thTags}, \texttt{thName}, \texttt{thModel}, \texttt{thGround} &
\(X\) & Section III.T: the five methods and their names, the model of
\(V^{(2)}_{\rm SF}\) of a method, and its ground state at \(T=0\)
iterated to \(10^{-13}\) \\
\texttt{i1Grid}, \texttt{thCoup}, \texttt{thMassN}, \texttt{thLam},
\texttt{thTc} & \(I^{(1)}_{\rm SF}(M,T)\), \(g^{X,{\rm S}}_{\rm SF}\),
\(g^{X,{\rm P}}_{\rm SF}\), \(M^{X,{\rm n}}_{\rm SF}(T)\),
\(\lambda^{\rm P}_{\rm SF}(T)\), \(T^{X,{\rm c}}_{\rm SF}\) & Section
III.T: the integral of the pair response on the radial nodes, the
effective couplings, the Dirac mass of the normal state from its scalar
self-consistency, the eigenvalue of the pair response of the normal
state, and the critical temperature \\
\texttt{thSF}, \texttt{thGapPhi}, \texttt{thPhiGap00} &
\(\Phi^X_{0,{\rm SF}}(T)\) & Section III.T: the superfluid and the
normal state of CJT at \(0.5\) and \(0.9\,T^{\rm CJT,c}_{\rm SF}\), the
order parameter from the gap equation alone, and its values for HS
\((0,0)\) \\
\texttt{thCurveTaus}, \texttt{thRes}, \texttt{thRustRes} & -- & Section
III.T: the reduced temperatures of the free-energy curves, and the
residuals of the Landau form and of the quartic form recomputed by the
notebook and of the Rust program \\
\texttt{thRadius}, \texttt{thCross}, \texttt{thShown},
\texttt{thLegend}, \texttt{thColors} & \(r^X_{\rm SF}(T)\) & Section
III.T: the radius of convergence with the Dirac mass of the normal state
computed by the notebook, the reduced temperature above which the order
parameter lies inside it, and the curves, the legend and the colors of
the figure \\
\end{longtable}

\textbf{The names of the code of Part R.} Part R starts from a clean
slate (Section 10), so its names do not collide with those above; each
of its Sections R.2 to R.4 starts from a clean slate again, and a name
may stand for different objects in different sections of Part R (as
\texttt{mEff}, which is \(M^{\rm eff}_{\rm SH}\) wherever it occurs).
The names that stand for objects of the physics:

\begin{longtable}[]{@{}
  >{\raggedright\arraybackslash}p{(\columnwidth - 4\tabcolsep) * \real{0.3333}}
  >{\raggedright\arraybackslash}p{(\columnwidth - 4\tabcolsep) * \real{0.3333}}
  >{\raggedright\arraybackslash}p{(\columnwidth - 4\tabcolsep) * \real{0.3333}}@{}}
\toprule\noalign{}
\begin{minipage}[b]{\linewidth}\raggedright
name in the code
\end{minipage} & \begin{minipage}[b]{\linewidth}\raggedright
printed form
\end{minipage} & \begin{minipage}[b]{\linewidth}\raggedright
meaning
\end{minipage} \\
\midrule\noalign{}
\endhead
\bottomrule\noalign{}
\endlastfoot
\texttt{rvKeep}, \texttt{rvClean}, \texttt{rvCheck}, \texttt{rvRust},
\texttt{rvPython}, \texttt{rvData}, \texttt{readTable}, \texttt{num},
\texttt{col} & -- & R.0: the clean slate, the check of Part R with its
counterparts, the tables of the Rust programs and the reading of their
fields \\
\texttt{sig{[}p{]}}, \texttt{sigbar{[}p{]}}, \texttt{gam{[}p{]}},
\texttt{gch}, \texttt{cC}, \texttt{bet} & \(\sigma^p\),
\(\bar\sigma^p\), \(\gamma^p_{\rm LL}\), \(\gamma^{\rm ch}_{\rm LL}\),
\(C_{\rm LL}\), \(\beta_{\rm LL}\) & R.1, R.2: the Pauli and the Dirac
matrices of the chiral basis \\
\texttt{d8}, \texttt{blk}, \texttt{dirac16} & -- & R.1: a Dirac matrix
on \(\Psi_{\rm LL}\), a block-diagonal matrix, a Dirac matrix on the
sixteen components \\
\texttt{T{[}a{]}}, \texttt{TR{[}a{]}}, \texttt{FF}, \texttt{YY},
\texttt{RR}, \texttt{QQ}, \texttt{gens}, \texttt{gensAll} &
\(T^{a,{\rm SM}}\), \(T^{a,{\rm R}}_{\rm SH}\),
\(F^{\rm tot}_{\rm SH}\), \(Y^{\rm SM}\), \(R_{\rm SH}\), \(Q^{\rm SM}\)
& R.1: the internal generators on the sixteen components, in the order
of the tables of the Rust program \\
\texttt{lb}, \texttt{sGen}, \texttt{l16} & \(u\), \(S(u)\) & R.1: the
six real generators of \(SL(2,\mathbb C)\) and their action on the
fields \\
\texttt{invOps}, \texttt{chargeOp}, \texttt{parityOp}, \texttt{spaces},
\texttt{cspaces} & \(L_X\) & R.1: the conditions of Lorentz invariance
of a bilinear, the action \(L_X\) of a generator and of parity on it,
and the invariant spaces of the twelve and of the four custodial
sectors \\
\texttt{kp}, \texttt{km} &
\(C_{\rm LL}\gamma^{\rm ch}_{\rm LL}\otimes\mathbf e^s_{\rm SH}\),
\(C_{\rm LL}\otimes\mathbf e^s_{\rm SH}\) & R.1: the pair matrices of
positive and of negative parity \\
\texttt{hg}, \texttt{hg2}, \texttt{bg} & \(G_X\) & R.1: the generators
on \(\Phi_{\rm SH}\), on \(\Phi'_{\rm SH}\) and on the bidoublet
\(\mathcal M_{\rm SH}\) \\
\texttt{channels}, \texttt{cchannels}, \texttt{stab}, \texttt{st8},
\texttt{m2}, \texttt{ev12} & -- & R.1: the 24 and the 8 channels with
their quantum numbers, the stabilizer of a condensate, that of
\(i\tau^2\), the mass matrix of \((W^1,W^2,W^3,B)\) and its
eigenvalues \\
\texttt{k3}, \texttt{kP{[}a,\ s{]}} & \(\mathsf K^3\),
\(\mathsf K^{\rm d}_{\rm SH}\), \(\mathsf K'^{\rm u}_{\rm SH}\) & R.2:
the triplet pair matrix of Parts II and III and the pair matrices of the
revision \\
\texttt{i1closed}, \texttt{i1}, \texttt{gcritSH}, \texttt{lamExact},
\texttt{lamS}, \texttt{lamTrip} & \(I^{(1)}_{\rm SF}\),
\(G^{\rm crit}_{\rm SH}\), \(\lambda^{\rm H}_{\rm SH}\),
\(\lambda^{\rm S}_{\rm SH}\), \(\lambda^{\rm T}_{\rm SH}\) & R.2, R.3:
the integral of the gap equation, the critical coupling and the
eigenvalues of the normal vacuum \\
\texttt{mLL}, \texttt{gSH}, \texttt{mEff}, \texttt{dAbs}, \texttt{d0},
\texttt{phi0} & \(m_{\rm LL}\), \(G_{\rm SH}\),
\(M^{\rm eff}_{\rm SH}\), \(\lvert D^{\rm d}_{\rm SH}\rvert\),
\(d^0_{\rm SH}\), \(\varphi^0_{\rm SH}\) & R.2 to R.4: the mass, the
pairing coupling \(504/5\), and the ground state of the pairing-only
methods \\
\texttt{chi6}, \texttt{chi4}, \texttt{zz}, \texttt{phi},
\texttt{vppOne}, \texttt{jj} & \(\chi^{ab}_{\rm SH}\), \(Z_{\rm SH}\),
\(\phi\), \(V^{\rm eff\,\prime\prime}_{\rm SH}(\varphi^0_{\rm SH})\),
\(J^{(3)}_{\rm SH}\) & R.3: the susceptibilities of the effective
Lagrangian, its normalization and the length of the vacuum, the
curvature of the potential of one pair, the integral of its
derivative \\
\texttt{lamAmp}, \texttt{lamPh}, \texttt{wFound}, \texttt{thr} &
\(\lambda^{\rm amp}_{\rm SH}\), \(\lambda^{\rm ph}_{\rm SH}\),
\(\omega^{\rm H}_{\rm SH}\), \(2M^{\rm eff}_{\rm SH}\) & R.3: the
eigenvalues of the response, the amplitude mode and the threshold \\
\texttt{r12Dirs}, \texttt{r12Tv}, \texttt{r12Cas}, \texttt{r12Broken} &
\(T^{a,{\rm SM}}+T^{a,{\rm R}}_{\rm SH}\) & R.3: the eight directions of
a fluctuation of the bidoublet, the custodial generators on it, their
Casimir operator and the directions of the broken generators (remark
R12) \\
\texttt{hh}, \texttt{rr}, \texttt{fN}, \texttt{rates} &
\(\mathsf H_{\rm SH}\), \(\mathcal R_{\rm SH}\),
\(F^{\rm tot}_{\rm SH}\) & R.4: the symbolic Hermitian matrices of the
conservation laws, the particle number on the Nambu space, the Josephson
rates \\
\texttt{wHann}, \texttt{ww}, \texttt{hw} & \(h^{\rm Hann}_{\rm SH}\),
\(\Gamma^{\rm win}_{\rm SH}\) & R.4: the transform of the Hann window
and the width of a pure cosine \\
\texttt{eps2}, \texttt{w2}, \texttt{dl0}, \texttt{vx}, \texttt{vy},
\texttt{vz}, \texttt{sol}, \texttt{dps} &
\(\pm\varepsilon_{\rm LL}(q_i)\), \(w_i/2\), \(\Delta_{\rm SH}\),
\(\mathbf n^\pm_{\rm SH}\), \(d_{\rm SH}(t)\) & R.4: the reduced model
of the pseudospins, its solution by \texttt{NDSolve} with the method
IDA, and the pair self-energy it gives \\
\end{longtable}

\hypertarget{errors-and-corrections-found-while-this-version-was-written}{%
\section{9. Errors and corrections found while this version was
written}\label{errors-and-corrections-found-while-this-version-was-written}}

\begin{enumerate}
\def\labelenumi{\arabic{enumi}.}
\item
  \emph{A comment that ended early (E1).} Two comments of the code of
  theorem D1 contained the characters \texttt{\^{}*)} (in
  \texttt{s(u)\^{}*} and \texttt{(psi\^{}+)\^{}*}), which end a comment
  of the Wolfram language; the first dry run stopped with a syntax
  error. The comments now write the complex conjugate as
  \texttt{conj(...)}.
\item
  \emph{The table of the characters of TeX commands (E2).} The reader of
  formulas took this table from the reader of symbols of the notebook
  \texttt{3+1super-loop.nb}; without it the reader would not know the
  Greek letters. The table is now part of \texttt{sf\_formula.wl}.
\item
  \emph{A name that resembled a forbidden word (E3).} The first
  evaluated notebook named the coupling of the pairing term with three
  letters that the scan of every commit rejects (the scan refuses a
  short list of names in the files of the project). The couplings are
  now named \texttt{gScal}, \texttt{gVec} and \texttt{gPair}, and the
  notebook was built and evaluated again.
\item
  \emph{Two checks with one identifier (E4).} The table of Section 6
  compares every check with its counterparts by their identifiers. It
  showed that the Python program had used the identifier E-10 for its
  symbolic derivation of the Yang-Mills equations, while E-10 of the
  Rust program is a different check (the self-energy of Hartree and Fock
  is the derivative of the energy). The Python check is now E-14, the
  Rust program has a new check E-14 of the same statement (exact
  derivatives of the polynomial Lagrangian density), and the notebook's
  check is WE-14. The manuscript (Section 5.2 and Section 17.5) states
  this.
\item
  \emph{A check without counterparts (E5).} The notebook's check of the
  conventions of the Standard Model had at first no counterpart in the
  two programs; the check SM-1 was added to both (Section 1 of the
  manuscript).
\item
  \emph{The printed time of the dry run (E6).} The first dry run printed
  the computing times as unevaluated number formats; the dry run now
  prints rounded numbers. This affects only the dry run.
\item
  \emph{A figure wider than the printed page (E7).} The first evaluation
  of the section of Part III placed the two panels of its figure side by
  side, 880 points wide; the printed copy cut the right panel. The
  panels now stand one above the other, each 460 points wide, and the
  notebook was built and evaluated again.
\item
  \emph{A warning of the numerical integration (E8).} The first dry run
  of Part II gave \texttt{NIntegrate} the two radii as machine numbers
  with the working precision of 30 digits, and the kernel warned that
  the integrand had less precision than requested; the radii are now
  given 30 digits (\texttt{SetPrecision}), and the warning is gone.
\item
  \emph{A badly conditioned system (E9).} The first dry run of Section
  III.T solved the cubic interpolation of the Taylor coefficients in the
  variable \(\Phi^2\) (values from \(1.1\times10^{-4}\) to
  \(1.7\times10^{-3}\)), whose matrix has the condition number
  \(1.2\times10^{10}\), and \texttt{LinearSolve} warned that its result
  might contain significant numerical errors (the coefficients agreed
  with those of the Rust program to \(1.4\times10^{-7}\) all the same).
  The interpolation now uses the scaled variable
  \(\Phi^2/\Phi^2_{\max}\), with the condition number 146, and the
  warning is gone.
\item
  \emph{A sentence that the notebook had outgrown (E10).} The
  introduction of the first section of Part III said that Part II, whose
  engine computes such curves, ``follows in the next versions''; Part II
  was added in the version before this one. The sentence now says what
  that chapter takes from the Rust program and that the next chapter
  uses the engine of Part II.
\item
  \emph{One identifier for two checks (E11).} When Section III.T named
  the counterparts of its checks, the Python check of the uniqueness of
  the critical temperature had the identifier YP3-T5, which by the rule
  of the identifiers belongs to the Python counterpart of the Rust check
  P3-T5 (the order parameter of Landau's theory). The Python check is
  now YP3-T9, the Rust program makes the same check as P3-T9, and the
  Python program has the new counterparts YP3-T5 to YP3-T8 of the Rust
  checks P3-T5 to P3-T8, so that every check of Section III.T has a
  counterpart in both programs.
\item
  \emph{One identifier for two checks, again (E12).} When Section III.G
  named the counterparts of its checks, the Python checks of the
  gradient coefficient and of the Higgs sector had been numbered YP3-G1
  to YP3-G4 and YP3-H1 to YP3-H3 in the order in which the Python
  program made them, so that six of them carried the identifiers of Rust
  checks of other statements (the error of item 11 again). The Python
  program now makes the counterpart of every Rust check under the
  identifier of the rule, and its earlier checks of statements without a
  Rust check (the slope in closed form, the two forms of the Ginzburg
  number, the quartic coefficient of CJT, the numbers of the Higgs
  sector, the row of the Standard Model) received Rust checks of the
  same statements, P3-G10 to P3-G12, P3-H3 and P3-H4. Found before the
  notebook was evaluated.
\item
  \emph{Two sentences that the notebook had outgrown (E13).} Section
  III.T.7 said that the gradient coefficient at the critical temperature
  ``is not computed yet'' and that computing it is the next step;
  Section III.G computes it, and the sentence now states its result. The
  answer of Section III.T.6 (``smaller and more physically reasonable
  near the critical temperature'') now says that it holds for the
  mean-field free energy of the methods, because Section III.G shows
  that the temperatures where it holds lie inside the critical region.
\item
  \emph{A variable of the same name, and a number joined to a text
  (E14).} The first dry run of Section III.G found two errors of its
  code. (a) The function of the product rule numbered the angles
  \(\varphi\) with the loop variable \texttt{k}, the name of the pair
  momentum in the radius of the lens; the loop variable of the function
  shadowed the momentum, the lens was built with the values of the angle
  index, and check WP3-G9 failed for all five methods. The loop variable
  is now \texttt{iphi}, and the radius of the lens receives the value of
  the momentum inserted (\texttt{With}). (b) The table of the Ginzburg
  numbers joined a number to a text, and the kernel issued the message
  \texttt{StringJoin::string} (which the evaluated notebook must not
  have); the number is now turned into a text with \texttt{ToString}.
  The second dry run passed all checks without a message.
\item
  \emph{A bound that no program had computed (E15).} The text of Section
  III.G.8 said that the masses of the gauge bosons are ``at most about a
  fifth of \(M^{X,{\rm H}}_{\rm SF}/2\) (the phase-space factor is above
  0.97)''. In KS and CJT the larger of the masses of \(W\) and \(Z\) is
  \(0.213\) of \(M^{X,{\rm H}}_{\rm SF}/2\), and the sentence named only
  the factor of the phase space of the decays into \(W^+W^-\) and
  \(ZZ\), not that of the sum over the polarizations. It was found on 21
  November 2026, when the numbers typed in the drafts of the article for
  Physical Review D were compared with the data; the same sentence stood
  in Section 14.6 of the manuscript and in the addendum A2 of
  \texttt{provenance/P4\_remarks.md}, which were corrected first.
  Section III.G.8 now explains the corrections at tree level, with the
  factor as a formula, and a new input cell (\texttt{gzVV}) computes for
  every method the larger ratio of the masses, the smaller factor and
  how much the masses of the effective Lagrangian exceed those of the
  linear response. The cell makes no check: its numbers come from the
  values of the Rust program that the checks WP3-H1 to WP3-H3 have
  already compared with the notebook's own.
\item
  \emph{The electric charge (C18 and C19 of the task record).} The
  notebook, like the manuscript, called the unbroken charge the electric
  charge and the massless gauge boson the photon, and Section III.G.8
  judged the superfluid against the Higgs boson of the Standard Model
  (its ratios against 0.5085 and 0.1293, the bound \(8\pi/3\) of the
  Higgs doublet, the masses and the critical temperature in GeV on the
  scale \(v=246.22\) GeV, a table with a row of the Standard Model, and
  the check WP3-H4 of that row). On 21 November 2026 the author stated
  that the spinor field, the DiracLike spinor, has electric charge 0 and
  no electroweak quantum numbers, and that
  \(SU(2)_{\rm v}\times U(1)^{\rm F}_{\rm LL}\) is a new gauge group,
  not the electroweak group. The comparison had no basis and is
  withdrawn: Section III.G.8 now answers separately whether the
  superfluid is the Higgs field of its own group and whether it can be
  that of the Standard Model, its table has no row of the Standard Model
  and no column in GeV, and the check WP3-H4 was removed with its
  counterparts P3-H4 and YP3-H4. Section I.1 states that the conventions
  of the Standard Model give names only; the texts of Part II name the
  massless gauge boson the zeugon, and a remark at the end of Part II
  states what the computation establishes about it, including that its
  charge conjugation is a gauge symmetry (Alice electrodynamics); the
  spinor field is called the DiracLike field. The data cells were
  written again from the data of the programs run again after the
  correction (their kept numbers are unchanged; the cell of Section
  III.G no longer holds the columns in GeV and the row of the Standard
  Model), and the notebook was built and evaluated again on 21 November
  2026: ALL PASS (181 of 181), 0 messages. Since this evaluation the
  counts of the builder are no longer typed in Section 4 but written
  from its report.
\item
  \emph{Part R (22 November 2026).} (1) The cell of remark R12
  multiplied a matrix halved, written
  \texttt{r12Tau{[}{[}a{]}{]}/2\ .\ m}; the product of matrices binds
  more strongly than the division, so that this was
  \texttt{r12Tau{[}{[}a{]}{]}/(2\ .\ m)}. The first dry run of Part R
  reported seven independent directions of the broken generators instead
  of four and did not find the charges of the triplet; the products are
  now in parentheses, and the cell gives four Goldstone directions, the
  singlet and the triplet. (2) The clean slate of Part R also cleared
  the two names with which the dry run measures its times, so that it
  printed no times; names that begin with \texttt{SFNB}, those of the
  tools that evaluate the notebook, are now kept. (3) Numbers formatted
  by \texttt{NumberForm} inside \texttt{Print} appeared in the output of
  the dry run, which is plain text, as \texttt{NumberForm{[}...{]}};
  they are now turned into strings first. (4) When the counterparts of
  the checks of Part R were worked out from the reports, it was found
  that the checking programs of the revision had said that their
  identifiers were those of the Rust programs; that was false for stages
  3, 4 and 6, and the programs were corrected first (Section 4.3 of the
  task record). (5) In the printed copy of the first evaluation of Part
  R two display formulas were too long for one line and broke inside a
  symbol, the pairing terms with the pair operators of Section R.0 and
  the two eigenvalues of remark R11 in Section R.3.5; each is now two
  display formulas, and the notebook was built and evaluated again. (1)
  to (3) were found by dry runs, (5) by looking at the printed copy.
\end{enumerate}

\hypertarget{part-r-the-revised-hypotheses-22-november-2026}{%
\section{10. Part R: the revised hypotheses (22 November
2026)}\label{part-r-the-revised-hypotheses-22-november-2026}}

\textbf{The task.} On 21 November 2026 the author asked to modify the
hypotheses so that the superfluid can be the Higgs field of the Standard
Model, with the quantum numbers of the Higgs doublet, and then to add a
second singlet of hypercharge \(-\tfrac34\). The revision was carried
out in stages, each with a Rust program, a Python program and a Wolfram
script; its stage 6 has three parts: the dynamics in time, the notebook,
and the article. The notebook is the second part, and the remarks made
during the revision, R8 to R12, go into it as into the manuscripts and
the document of the remarks.

\textbf{What Part R holds.} Section R.0 states the fields, their
charges, the symmetries and the interaction of the second version,
explains how its checks are named, and defines its tools: the clean
slate, the check \texttt{rvCheck} with the counterparts of every check,
and the reading of the tables of the closed cell of data. Section R.1
builds the generators on the sixteen components as exact matrices and
derives the quantum numbers (WQN-1 to WQN-12), with remark R8 and the
products of the charges on which it rests. Section R.2 derives the
formulas of the ground state (WSF-2a to WSF-5), with remark R9 and the
charges of the quasi-particles of the second version, and remark R10
with the surface integral of the sharp cutoff in closed form, against
the eigenvalue of the Rust program. Section R.3 derives the custodial
algebra of the effective Lagrangian, the mass of its amplitude mode, the
closed forms of HS (0, 0), the scale in GeV, and the amplitude mode at
\(d^0_{\rm SH}\) (WH-1 to WH-5), with remark R11 and remark R12, whose
cell finds the four Goldstone directions, the massive singlet and the
massive custodial triplet of the bidoublet and the charges \(+1\), 0,
\(-1\) of the triplet. Section R.4 derives the conservation laws, the
relation of Josephson and the width of the Hann window, compares the
peaks of the Rust program, and integrates the 48 pseudospins of the
reduced model with \texttt{NDSolve} and the method IDA (WD-1 to WD-5).

\textbf{How it is made.} The code of the 28 checks is that of the four
Wolfram scripts of the revision, copied by a program and not retyped,
without the lines that read the arguments, define the report and write
it; the scripts' own check \texttt{check{[}id,\ text,\ ok,\ detail{]}}
became \texttt{rvCheck} with the same arguments, which records the check
in the register of Chapter 0 together with its counterparts. Every table
that a script read from a file is in the closed cell of data. Each of
the four scripts ran in a kernel of its own; in the notebook each
section of Part R starts from a clean slate that clears every name of
the global context except those of the register, of the tools of Part R
and of the tools that evaluate the notebook, so that a short name of the
algebra cannot carry a value from an earlier part or section.

\textbf{The evaluation.} The whole notebook, Parts I to III and Part R,
was built and evaluated in the front end on 22 November 2026 (Section
5): ALL PASS (181 of 181), 0 messages.

\hypertarget{what-follows}{%
\section{11. What follows}\label{what-follows}}

The rest of Part III and Part IV of the notebook for the earlier
hypotheses (the effective Lagrangians, the masses of the gauge bosons in
every vacuum, the surviving symmetry, the dynamics with IDA inside
\texttt{NDSolve}, the figures and the videos) would be added in the same
way, with the checks of the Python programs \texttt{sf\_part3.py} and
\texttt{sf\_part4.py}. Under the revised hypotheses, Part R grows with
every further stage of the revision, and this document with it.

\end{document}
